\exercisesection{4}

\begin{enumerate}
\def\labelenumi{\arabic{enumi}.}
\tightlist
\item
  \begin{enumerate}
  \def\labelenumii{(\alph{enumii})}
  \tightlist
  \item
    设 G 为全序群，M 为 \(G_{+}\) 中不含 0 的优势集。证明：存在 G 中与 M 不相交的最大孤立子群 H。
  \end{enumerate}
\end{enumerate}

\begin{enumerate}
\def\labelenumi{(\alph{enumi})}
\setcounter{enumi}{1}
\item
  设 A 为赋值环，v 为与 A 相伴的赋值。A 的理想 \(p \neq 0\) 为素理想的充要条件是：p 对应于 v 的全序取值群 G 中的一个优势集 M，并且若 H 为 G 中与 M 不相交的最大孤立子群，则 M 是 \(H_{+}\) 在 \(G_{+}\) 中的补集。
\item
  沿用 (b) 的记号，A 的理想 q 为 p-准素理想的充要条件是满足下列条件之一：或者 \(H = \{0\}\)（此时 p 为极大理想），或者 q = p。
\item
  对于 \(\mathfrak{a}\) 这一 \(\mathbf{A}\) 的理想，若它既不等于 0 也不等于 \(\mathbf{A}\)，则满足 \(\mathfrak{a}^2 = \mathfrak{a}\) 的充要条件是：\(\mathfrak{a}\) 对应于一个优势集 \(\mathbf{M}\)，且若 \(\mathbf{H}\) 是 \(\mathbf{G}\) 中与 \(\mathbf{M}\) 不相交的最大孤立子群，则 \(\mathbf{M}'\) 是 \(\mathbf{M}\) 在全序群 \(\mathbf{G}' = \mathbf{G} / \mathbf{H}\) 中的像，且没有最小元；此时 \(\mathfrak{a}\) 是素理想。
\end{enumerate}

\begin{enumerate}
\def\labelenumi{\arabic{enumi}.}
\setcounter{enumi}{1}
\tightlist
\item
  设 G 为全序群。
\end{enumerate}

\begin{enumerate}
\def\labelenumi{(\alph{enumi})}
\item
  对于 G 的每个元素 a \textgreater{} 0，令 \(\mathrm{H}(a)\) 为包含 a 的最小孤立子群，令 \(\mathrm{H}^{-}(a)\) 为不包含 a 的最大孤立子群；证明 \(\mathrm{H}(a)/\mathrm{H}^{-}(a)\) 同构于 R 的一个子群。
\item
  反之，若 H 是 G 的孤立子群，且包含于 H 的 G 的孤立子群所成集合中存在关于包含关系的最大元 \(H'\) \(\neq H\)，则 \(\mathrm{H} = \mathrm{H}(a)\) 且 \(\mathrm{H}' = \mathrm{H}^{-}(a)\)，对所有 \(a \in H_{+} \cap CH_{+}'\) 均成立。这样的孤立子群 H 称为主孤立子群，\(H'\) 称为它的前驱。
\item
  若 \(\mathbf{H}_1, \mathbf{H}_2\) 是 \(\mathbf{G}\) 的两个不同孤立子群，且 \(\mathbf{H}_1 \subset \mathbf{H}_2\)，证明存在 \(\mathbf{H}, \mathbf{H}'\) 的两个孤立子群 \(\mathbf{G}\)，使得 \(\mathbf{H}_1 \subset \mathbf{H}' \subset \mathbf{H} \subset \mathbf{H}_2\)，并且 \(\mathbf{H} / \mathbf{H}'\) 同构于一个非零的 \(\mathbf{R}\) 子群。
\item
  令 \(\Sigma\) 为 \(G\) 的孤立子群所成的集合（按包含关系全序），令 \(\Theta \subset \Sigma\) 为主孤立子群所成的集合。证明 \(\Sigma\) 同构于 \(\Theta\) 的完备化（集合论，第三章 § 1，习题 15）。
\end{enumerate}

\begin{enumerate}
\def\labelenumi{\arabic{enumi}.}
\setcounter{enumi}{2}
\tightlist
\item
  \begin{enumerate}
  \def\labelenumii{(\alph{enumii})}
  \tightlist
  \item
    令 \(\Theta\) 为全序集，并对每个 \(s \in \Theta\) 令 \(\mathbf{A}_s\) 为不退化为 0 的 \(\mathbf{R}\) 的子群。令
  \end{enumerate}
\end{enumerate}

\[
\Gamma (\Theta , (A _ {s}) _ {s \in \Theta}) = G
\]

是由支集为 \(\biguplus_{s\in\Phi}\mathbf{A}_{s}\) 的良序子集的函数组成的乘积 \(f\) 的子群。我们在 \(\Theta\) 上定义一个与其群结构相容的序：令元素 \(\mathbf{G}\) 所成的集合 \(\mathbf{G}_{+}\) 由函数 0 以及满足下述条件的函数 \(\geqslant0\) 组成：f 在其支集的最小元 \(f\neq0\) 处满足 \(f(\theta)>0\)。证明 \(\theta\) 是全序群，并且 \(f\) 按关系 \(\mathbf{G}\) 排序后典范同构于 \(\Theta\) 的主孤立子群所成的集合；此外，在此同构下若 \(\mathbf{G}\) 对应于 \(\supset\)，则 \(\mathrm{H}(a)\) 同构于 \(s\in\Theta\)。\(\mathrm{H}(a)/\mathrm{H}^{-}(a)\)\(\mathbf{A}_{s}\)

\begin{enumerate}
\def\labelenumi{(\alph{enumi})}
\setcounter{enumi}{1}
\tightlist
\item
  考虑全序群 \(G'\)（按字典序排序）的子群 \(Q \times Q\)，它由元素 \((p_{n}^{-1}, np_{n}^{-1})\) 生成，其中 \((p_{n})\) 是严格递增的素数序列。证明 G' 中既不同于 0 也不同于 \(H'\) 的唯一孤立子群 \(G'\) 是 \(G'\)；但 G' 不同构于按字典序排序的乘积 \(\{0\} \times Z\)\(G'\)\(H' \times (G'/H')\) 的任何子群。
\end{enumerate}

\begin{enumerate}
\def\labelenumi{\arabic{enumi}.}
\setcounter{enumi}{3}
\tightlist
\item
  设 \(A\) 为非域的赋值环。
\end{enumerate}

\begin{enumerate}
\def\labelenumi{(\alph{enumi})}
\item
  证明 A 为离散赋值环的充要条件是 A 的每个素理想都是主理想。
\item
  证明：若 A 是高度 1 的赋值的环，则 A 的分式域是有限生成的 A-代数。
\end{enumerate}

* 5. 设 K 为域，M 为 K 中非域子环所成的集合。证明 M 的极大元是其分式域 L 的一个高度 1 的赋值环，并且 K 是 L 的代数扩张。下列条件等价：

(α) M 非空。

(β) M 有极大元。

(γ) K 上存在高度 1 的赋值

(δ) K 不是有限素域的代数扩张（参见 § 8，第 1 节）。*

¶ 6. (a) 设 S 为没有孤立点的超 Stonean 紧空间（积分论，第五章 § 5，习题 14）。令 \(\mathcal{C}_{0}(S)\) 为 S 上取实值（有限或非有限）、在 S 上连续且满足集合 \(\overline{f}(-\infty)\) 与 \(\overline{f}(+\infty)\) 在 S 中处处不稠密的函数所成的向量空间（积分论，第二章 §1，习题 13 (g)）。证明：对 S 上每个测度 \(\mu > 0\)，\(\mathcal{C}_{0}(S)\) 中存在上积分为无穷的函数 f（注意，对于 S 的每个处处不稠密闭子集 N，都存在 \(\mathcal{C}_{0}(S)\) 中在 N 上等于 \(+\infty\) 的函数 f；然后证明只需考虑测度 \(\mu\) 为正规测度的情形，并将 f 适当地定义为 \(\mathcal{C}_{0}(S)\) 中一列函数的上界）。

\begin{enumerate}
\def\labelenumi{(\alph{enumi})}
\setcounter{enumi}{1}
\item
  令 \(\mathbf{G}\) 为 \(\mathcal{C}_0(S)\) 的有序子群，它由 \(\mathcal{C}_0(S)\) 中取值于 \(\mathbf{Z}\) 或 \(\pm \infty\) 的函数组成。证明 \(\mathbf{G}\) 是完备格，并且（作为有序群）不同构于乘积群 \(\mathbf{R}^{\mathbf{l}}\) 的任何子群。
\item
  令 K 为域， 为以空间 S 为指标集的一族不定元（\(A_{0} = K[[X_{s}]]_{s \in S}\)）的系数在 K 中的形式幂级数环（代数，第 IV 章 § 5，习题 1）。令 b 为满足下述性质的形式幂级数 \(X_{s}\) 所成的集合：存在 S 的处处不稠密子集 N，使得 P 的每个系数 \(P \in A_{0}\) 的单项式都含有某个满足 \(N \subset S\)\(\neq 0\) 的 。证明 b 是 \(X_{s}\) 的理想，并且环 \(s \in N\)\(A_{0}\) 是整环且完全整闭（证明后一点时应用第五章 § 1，第 4 节的命题 14，并仿照第五章 § 1，习题 10 论证；还要使用 S 中每个第一类集都是处处不稠密的事实）。\(A = A_{0}/b\) 是整环且完全整闭（证明后一点时应用第五章 § 1，第 4 节的命题 14，并仿照第五章 § 1，习题 10 论证；还要使用 S 中每个第一类集都是处处不稠密的事实）。
\item
  令  为群 \(f\) 中满足  的元素，令 N 为 \(\geqslant 0\) 的点所成的处处不稠密集合。令 \(G\)\(N\)\(f(s) = \pm \infty\) 为形式幂级数 \(p_f\) 在 A 中的像；映射 \(A\)\(\prod_{s \notin N} (1 + X_s)^{f(s)}\) 是从加法幺半群 \(f \mapsto p_f\) 到 A 的一个乘法子幺半群的同构。推出 A 不可能是其分式域的高度 1 的一族赋值环的交（证明这将蕴含 G 同构于乘积 \(G_+\) 的一个子群）。（参见习题 8、9。）\(A\)\(A\)\(G\)\(R^1\) 的一个子群）。（参见习题 8、9。）
\end{enumerate}

¶ 7. 若局部整环 A 不是域，且 A 除 (0) 与极大理想 m 外没有别的素理想，则称 A 的维数为 1。对 A 的每个理想 a，令 A: a 表示 A 的分式域中 a 在 A 中的传递子所成的 A-子模（它总是包含 A）。以下均假设 A 的维数为 1。

\begin{enumerate}
\def\labelenumi{(\alph{enumi})}
\item
  证明：若 A 是强 Lasker 环（第四章 § 2，习题 28），则 A: m ≠ A。（反设相反，注意对每个整数 r \(^{r}\) 0 都有 A: m \textgreater{} = A；另一方面，A 中包含于 m 的每个理想 ≠ 0 都是 m-准素的，因而特别地，每个主理想 Ax ⊆ m（其中 x ≠ 0）都是 m-准素的；最后注意 h ≠ k 时 m \(^{h}\) ≠ m \(^{k}\)（第四章 § 2，习题 29 (d)）。）
\item
  证明：对于维数为 1 的强 Lasker 局部整环 A，下列性质等价：
\end{enumerate}

\((\alpha)\) A 是完全整闭的。

(β) 极大理想 m 是主理想。

(γ) 每个 m-准素理想都形如 \(m^{k}\)。

(δ) 理想 m 是可逆的（第二章 § 5，第 6 节），这等价于 m. (A: m) = A。

(ε) A 是离散赋值环。

（为证明  蕴含 \((\alpha)\)，使用 (a)，并注意若对所有 \((\delta)\) 都有 ，则 \(\mathfrak{m}.\left(\mathbf{A}:\mathfrak{m}\right)^k = \mathfrak{m}\) 不是完全整闭的（参见第七章 §1，习题 4）。为证明 \(k > 0\)\(\mathbf{A}\) 蕴含 \((\delta)\)，注意对于每个包含于 \((\gamma)\) 的理想 ，都可写成 \(\mathfrak{a} \neq 0\)，其中 \(\mathfrak{m}\)\(\mathfrak{a} = \mathfrak{ma}_1\)，然后使用第四章 §2，习题 29 (d)。最后，为证明 \(\mathfrak{a}_1 \neq 0\) 蕴含 \((\gamma)\)，注意 \((\beta)\)，且 \(\mathfrak{m} \neq \mathfrak{m}^2\) 蕴含 \((\gamma)\) 是一个 1 维的 \(\mathfrak{m} / \mathfrak{m}^2\)-向量空间。）\((\mathbf{A} / \mathfrak{m})\)-向量空间。）

\begin{enumerate}
\def\labelenumi{(\alph{enumi})}
\setcounter{enumi}{2}
\item
  证明：用第三章 § 3，习题 15 (b) 的记号，局部环 \(\mathbf{A}_{\mathfrak{m}}\) 是维数为 1 的 Noether 整环，但不是整闭的。
\item
  设 K 为代数闭域，A 为域 \(\mathbf{K}(\mathbf{X},\mathbf{Y})\) 的子环；该域是 K 上两个不定元的有理函数域，而 A 由满足下述条件的元素 \(f\in\mathbf{K}(\mathbf{X},\mathbf{Y})\) 组成：可将 X 代入 f 中的 0，且 \(f(0,\mathbf{Y})\) 属于 K。证明 A 是强 Lasker 的整闭局部整环，维数为 1，但不是完全整闭的。（若 B 是多项式环 \(\mathbf{K}[\mathbf{X},\mathbf{Y}]\)，p 是 B 的素理想 BX，注意 A 是 \(B_{p}\) 中等于 \(K+pB_{p}\) 的子环。）(*)
\end{enumerate}

¶ 8. 设 A 为维数 1 的局部整环（习题 7），K 为其分式域。证明 A 是（K 的）高度 1 的赋值环的交的充要条件是 A 完全整闭。依次证明：

\begin{enumerate}
\def\labelenumi{(\alph{enumi})}
\item
  含有 A 以及 A 的极大理想 m 中一个元素 \(\neq0\) 的逆元的 K 的子环等于 K（使用 A 中每个既不等于 0 也不等于 A 的理想都是 m-准素理想这一事实）。
\item
  若 \(\mathbf{B} \supset \mathbf{A}\) 是不同于 \(\mathbf{K}\) 的 \(\mathbf{K}\) 的子环，则存在高度 1 的赋值环 \(\mathbf{V}\)，使得 \(\mathbf{B} \subset \mathbf{V}\)（使用 (a)）。
\item
  令 \(z \in \mathbf{K} - \mathbf{A}\)；仿照习题 7 (b) 论证 \(zm \subset m\) 不可能。推出存在高度 1 的赋值环 \(V\)，使得 \(A \subset V\) 且 \(z \notin V\)。（若 \(zm \subset A\)，使用 (a)；否则存在 \(x \in m\)，使得
\end{enumerate}

\[
x y = z \in \mathrm{K} - \mathrm{A};
\]

注意 \(z \notin A[z^{-1}]\)，并使用 (b) 证明 V 的存在性。）

¶ 9. 设 A 为 Lasker 完全整闭整环（第四章 §2，习题 23）。进一步假设，对 A 中所有 ，与 A/Ax 弱伴随的素理想 \(x \neq 0\)（第四章 § 1，习题 17）全都具有高度 1，也就是说对所有 i，\(\mathfrak{p}_i\) 不含除自身与 0 以外的素理想（参见第七章 § 1，第 6 节）。证明在这些条件下 A 是高度 1 的赋值环的交。（先使用第四章 § 1 的习题 17 (i)，证明 A 是所有环  的交，其中 p 遍历高度 1 的素理想所成集合。然后使用第四章 § 2，习题 23 的条件（），证明这些局部环  完全整闭。最后应用习题 8。）当 A 是每个理想都具有单个准素分解的 Lasker 完全整闭整环时，上述假设特别得到满足（第四章 § 2，习题 26）。\(p_{i}\)\(A_{p}\)\(A_{p}\)\(LA_{I}\) 完全整闭。最后应用习题 8。）当 A 是每个理想都具有单个准素分解的 Lasker 完全整闭整环时，上述假设特别得到满足（第四章 § 2，习题 26）。

¶ 10. 设 A 为一个主理想集合按包含关系全序的环。

\begin{enumerate}
\def\labelenumi{(\alph{enumi})}
\item
  证明 \(\mathbf{A}\) 是局部环，且其中的幂零根 \(\mathfrak{N}\) 是素理想。证明或者 \(\mathfrak{N}^2 = \mathfrak{N}\)，或者 \(\mathfrak{N}\) 是幂零的；环 \(\mathbf{A} / \mathfrak{N}\) 是赋值环。
\item
  若 \(\mathcal{L}\) 是 A 的主理想所成的全序集，证明 A 的理想所成的集合按包含关系全序，并同构于集合 \(\mathcal{L}\) 的完备化（集合论，第三章 § 1，习题 15）。
\item
  假设 ；于是可将 \(\mathfrak{N}^2 = 0\) 看作赋值环 \(\mathfrak{N}\) 上的模；令 K 为 V 的分式域，令 \(V = A / \mathfrak{N}\) 为 K 上与 V 对应的赋值 v 的阶群；证明 \(K\)\(V\)\(\Gamma\) 中存在两个优势集 \(v\)\(K\)\(V\)\(\Gamma\)，使得 \(M \subset M'\) 作为 V-模同构于 \(\mathfrak{N}\)（§3，第 4 节的记号）。\(V\)\(\mathfrak{a}(M') / \mathfrak{a}(M)\)（§3，第 4 节的记号）。
\item
  反之，给定分式域为 K、阶群为  的赋值环 V，以及 \(\Gamma\) 中的两个优势集 ，令 \(M \subset M'\)，并在乘法群 \(\Gamma\)\(Q = a(M') / a(M)\) 上规定乘法为\(A = V \times Q\) 上规定乘法为
\end{enumerate}

\[
(z, t) (z ^ {\prime}, t ^ {\prime}) = (z z ^ {\prime}, z t ^ {\prime} + z ^ {\prime} t);
\]

证明在 A 中主理想所成集合按包含关系全序；有序对 \((0, t)\)（其中 \(t \in Q\)）所成的集合 N 是 A 的幂零根，\(N^{2} = 0\)，且 A/N 同构于 V。

\begin{enumerate}
\def\labelenumi{(\alph{enumi})}
\setcounter{enumi}{4}
\item
  对每个素数 ，环 \(p\) 的主理想所成集合按包含关系全序，且 A 的幂零根 \(A = \mathbf{Z} / p^2\mathbf{Z}\) 满足 \(A\)\(\mathfrak{N}\)；但证明环 A 不同构于由 (d) 的方法（从 \(A\)\(\mathfrak{N}^2 = 0\) 与 \(A\)\(V = A / \mathfrak{N}\) 出发）构造的环（注意 A 不含同构于 \(Q = \mathfrak{N}\) 的子环）。\(A\)\(Z / p\mathbf{Z}\) 的子环）。
\item
  证明对于每个素数 ，群 \(p\) 关于素域 \(\mathbf{A}\)\(\mathbf{U}_p\) 的代数 \(\mathbf{F}_p\)（代数，第 VII 章 §2，习题 3）是一个环，其中主理想所成集合按包含关系全序，且幂零根  满足 \(\mathfrak{N}\)，但 \(\mathfrak{N}^2 = \mathfrak{N}\) 不同构于赋值环关于某个理想的商环。\(\mathbf{A}\) 不同构于赋值环关于某个理想的商环。
\end{enumerate}

\begin{enumerate}
\def\labelenumi{\arabic{enumi}.}
\setcounter{enumi}{10}
\tightlist
\item
  设 K 为带有高度 1 赋值 v 的域。
\end{enumerate}

\begin{enumerate}
\def\labelenumi{(\alph{enumi})}
\tightlist
\item
  令  为 \(\mathbf{P}(\mathbf{X}) = a_0\mathbf{X}^n +a_1\mathbf{X}^{n - 1} + \dots +a_n\) 中次数为 \(\mathrm{K[X]}\) 且满足 \(n > 1\) 的多项式；证明存在区间 \(a_{n} \neq 0\) 中严格递增的整数序列 ，使得：(1) \((i_{k})_{0 \leqslant k \leqslant r}\)，\([0, n]\)\(i_{0} = 0\)；(2) 对 \(i_{r} = n\)， 有限；(3) 对每个不同于各 \(v(a_{i_{k}})\) 的指标 ，若 \(0 \leqslant k \leqslant r\)\(j\) 且 \(0 \leqslant j \leqslant n\) 有限，则该点\(i_{k}\)\(v(a_{j})\) 有限，则该点
\end{enumerate}

\[
(j, v (a _ {j})) \in \mathbf {R} ^ {2}
\]

位于经过点 \((i_{k}, v(a_{i_{k}}))\) 与 \((i_{k+1}, v(a_{i_{k+1}}))\) 的直线之上；若 \(j < i_{k}\) 或 \(j > i_{k+1}\)，则严格位于该直线之上。连接点 \((i_{k}, v(a_{i_{k}}))\) 与 \((i_{k+1}, v(a_{i_{k+1}}))\) 的线段的并称为 P 的 Newton 多边形，上述线段称为边，点 \((i_{k}, v(a_{i_{k}}))\) 称为多边形的顶点。

\begin{enumerate}
\def\labelenumi{(\alph{enumi})}
\setcounter{enumi}{1}
\item
  假设 P 的所有零点都属于 K。证明 P 的所有零点具有相同赋值的充要条件是 r = 1（换言之，Newton 多边形退化为单条边）。（为证明条件充分，考虑乘积 \(P_{1}P_{2}\) 的 Newton 多边形，其中 \(P_{1}\) 的所有零点在 v 的环中均可逆，而 \(P_{2}\) 的所有零点都属于 v 的理想。）
\item
  假设 \(\mathbf{P}\) 的所有零点都属于 \(\mathbf{K}\)；作出 \(\mathbf{P}\) 的 Newton 多边形并令
\end{enumerate}

\[
\rho_ {k} = i _ {k + 1} - i _ {k}, \quad \sigma_ {k} = (v (a _ {i _ {k + 1}}) - v (a _ {i _ {k}})) / \rho_ {k}.
\]

证明：对 ，P 恰有 \(0 \leqslant k \leqslant r - 1\) 个零点（按重数计），其赋值都等于 \(\rho_k\)（使用 (b)，并对 r 作归纳）。\(\sigma_k\)\(r\)（使用 (b)，并对 r 作归纳）。

\begin{enumerate}
\def\labelenumi{(\alph{enumi})}
\setcounter{enumi}{3}
\tightlist
\item
  推广到任意赋值 v 的情形（将 v 的阶群 \(\Gamma\) 嵌入向量 Q-空间 \(\Gamma_{(Q)}\) 中，后者具有自然的全序群结构）。
\end{enumerate}

\exercisesection{5}

¶ 1. 设 A 为整环，K 为其分式域，T 为 A 上的线性拓扑（第三章 § 2，习题 21）。

\begin{enumerate}
\def\labelenumi{(\alph{enumi})}
\item
  要使 \(\mathcal{T}\) 下 0 的邻域构成某个与 K 的环结构相容的拓扑 \(\mathcal{T}_{\mathrm{K}}\) 在 0 处的基本邻域系，其充要条件是 \(\mathcal{T}\) 为拓扑 \(\mathcal{T}_{u}(\mathrm{A})\)（第三章 §2，习题 24）；此时相对于 \(\mathcal{T}_{\mathrm{K}}\)，A 是有界子集，而 \(\mathcal{T}_{\mathrm{K}}\) 是局部有界的 Hausdorff 拓扑（一般拓扑，第三章 §6，习题 12 与 20 (e)）。K 关于 \(\mathcal{T}_{\mathrm{K}}\) 完备（分别线性紧、严格线性紧（第三章 §2，习题 21））的充要条件是 A 关于 \(\mathcal{T}\) 具有相应性质。
\item
  拓扑 （其中 \(\mathcal{T}_{\mathbf{K}}\)）与 \(\mathcal{T} = \mathcal{T}_u(\mathbf{A})\) 的域结构相容的充要条件是 \(\mathbf{K}\) 的 Jacobson 根  满足 \(\Re(\mathbf{A})\)。\(\mathbf{A}\)\(\neq 0\)。
\end{enumerate}

¶ 2. 设 K 为（交换）域，T 为 K 上与其环结构相容的非离散 Hausdorff 拓扑。拓扑 T 由 K 上的赋值或绝对值定义的充要条件是 T 局部逆有界（一般拓扑，第三章 § 6，习题 22）。若 K 中存在拓扑幂零元，使用一般拓扑第三章 § 6 的习题 22 (d) 以及一般拓扑第九章 § 3 的习题 13。反之则使用一般拓扑第三章 § 6 的习题 22 (f)。

¶ 3. 设 A 为 Noether 整环，K 为其分式域，\(\mathcal{T}_{u}\) 为拓扑 \(\mathcal{T}_{u}(\mathrm{A})\)，\(\mathcal{T}_{\mathrm{K}}\) 为 K 上的相应拓扑（习题 1）。

\begin{enumerate}
\def\labelenumi{(\alph{enumi})}
\item
  若 A 是 Jacobson 根为 \(\mathfrak{r} \neq 0\) 的 Zariski 环，且 A 关于 r-进拓扑完备，证明 A 关于拓扑 \(\mathcal{T}_u\) 完备（一般拓扑，第三章 § 3，第 5 节，命题 9 的推论 2）。
\item
  假设  非离散且由 \(\mathcal{T}_{\mathbf{K}}\) 上的赋值 \(v\)\(\mathbf{K}\) 定义。证明 v 是离散赋值，且 \(v\) 是 v 的环 \(\mathbf{A}\) 的开子环；并证明逆命题。（利用蕴含 \(\mathbf{V}\) 的假设以及 §4 第 1 节的命题 1，可设 \(v\)\(\mathbf{A} \neq \mathbf{K}\)。利用 \(\mathbf{A} \subset \mathbf{V}\) 在 \(\mathbf{A}\) 下是开的事实，证明 \(\mathcal{T}_{\mathbf{K}}\) 是有限生成的 A-模，并使用 §3，第 5 节的命题 9 得出结论。对逆命题，注意 \(\mathbf{V}\) 是局部环，其中极大理想 \(\mathbf{A}\) 与 (0) 是唯一的素理想，因而其中每个理想 \(m\) 都是 \(a \neq 0\)-准素理想。此时 \(m\) 的整闭包为 \(\mathbf{A}\)。给出 \(\mathbf{V}\) 的例子（令 \(\mathbf{A} \neq \mathbf{V}\) 为形式幂级数环 \(\mathbf{V}\)，其中 k 为域）。\(k[[T]]\)\(k\)，其中 k 为域）。
\item
  推出一个非离散完备拓扑域的例子，其拓扑不是局部逆有界的（使用 (b) 和习题 2）。
\end{enumerate}

\begin{enumerate}
\def\labelenumi{\arabic{enumi}.}
\setcounter{enumi}{3}
\tightlist
\item
  设  为域 K 上的赋值，A 为 v 的环，m 为其极大理想。\(v\)\(K\)\(v\)\(m\) 为域 K 上的赋值，A 为 v 的环，m 为其极大理想。
\end{enumerate}

\begin{enumerate}
\def\labelenumi{(\alph{enumi})}
\item
  A 上由 v 定义的拓扑与 m-进拓扑相同的充要条件是 A 为域或离散赋值环。
\item
  由  定义的拓扑使 A 成为严格线性紧环（第三章 §2，习题 21）的充要条件是 A 为域或离散赋值环（使用 (a) 及第三章 §2，习题 22 (a)）。\(v\)\(A\)\(A\) 定义的拓扑使 A 成为严格线性紧环（第三章 §2，习题 21）的充要条件是 A 为域或离散赋值环（使用 (a) 及第三章 §2，习题 22 (a)）。
\end{enumerate}

\begin{enumerate}
\def\labelenumi{\arabic{enumi}.}
\setcounter{enumi}{4}
\tightlist
\item
  \begin{enumerate}
  \def\labelenumii{(\alph{enumii})}
  \tightlist
  \item
    设 K 为域，v 为 K 上的赋值，\(\Gamma\) 为 v 的阶群。K 关于拓扑 \(T_{v}\) 线性紧（第三章 §2，习题 15）的充要条件是：对 \(\Gamma\) 的每个良序子集 B 以及 K 的每个元素族 \((a_{\beta})_{\beta\in\mathcal{B}}\)，若对 \(\lambda<\mu<v\)，
  \end{enumerate}
\end{enumerate}

\[
v (a _ {\lambda} - a _ {\mu}) <   v (a _ {\mu} - a _ {\nu}),
\]

则存在元素 ，使得对满足 \(a \in \mathbf{K}\) 的每个指标有序对 ，都有 。（使用集合论第三章 § 2 的习题 4，\(v(a - a_{\lambda}) = v(a_{\lambda} - a_{\mu})\)\(\lambda, \mu\)\(\mu > \lambda\)。（使用集合论第三章 § 2 的习题 4，

第三章 § 2）。若此条件成立，则赋值 v 的环 A 关于离散拓扑也线性紧。\(v\)第三章 § 2）。若此条件成立，则赋值 v 的环 A 关于离散拓扑也线性紧。

\begin{enumerate}
\def\labelenumi{(\alph{enumi})}
\setcounter{enumi}{1}
\tightlist
\item
  证明 §3，习题 2 中定义的域  关于 \(S(\Gamma, k) = K\) 线性紧。令 \(T_{v}\) 为有理数的全序群 Q，并考虑 K 的子环 \(\Gamma\)，它由满足下述条件的 \(K_{0}\) 组成：使 \(x = (x_{\alpha})\) 的  所成集合是有限集，或是趋于 \(\alpha \in Q\) 的严格递增序列的点集。证明 \(x_{\alpha} \neq 0\)\(+\infty\) 是域，关于 K 上的典范赋值 v 所诱导的赋值 \(K_{0}\) 完备，并且 \(v_{0}\) 与 v 具有相同的取值群和剩余域，但 \(v_{0}\) 不是线性紧的（参见 §10，习题 2）。\(K_{0}\) 不是线性紧的（参见 §10，习题 2）。
\end{enumerate}

¶ 6. (a) 设 A 为赋值环，M 为 A-模，M' 为 M 的子模。证明 M' 为 M 的纯子模（第一章 § 2，习题 24）的充要条件是：对所有 \(\alpha \in A\)，M' ∩ (\(\alpha\) M) = \(\alpha\) M'。

\begin{enumerate}
\def\labelenumi{(\alph{enumi})}
\setcounter{enumi}{1}
\item
  设  是 \(\mathbf{M}'\) 的纯子模，并且若在 \(\mathbf{M}\) 上赋予这样的拓扑：\(\mathbf{M}'\)（其中 \(\alpha\mathbf{M}'\)）构成 0 的基本邻域系，则 \(\alpha \in \mathbf{A}, \alpha \neq 0\) 线性紧。证明对每个元素 \(0\)\(\mathbf{M}'\)，存在 \(x \in \mathbf{M}\)，使得 \(y' \in \mathbf{M}'\) 具有如下性质：对所有满足 \(x' = x + y'\) 的 ，都有 \(\alpha \in \mathbf{A}\)。（对于每个满足  的 \(x + \mathbf{M}' \in \alpha(\mathbf{M} / \mathbf{M}')\)，考虑 \(x' \in \alpha\mathbf{M}\)\(\alpha \in \mathbf{A}\) 的子集 ，它由满足 \(x + \mathbf{M}' \in \alpha(\mathbf{M} / \mathbf{M}')\)\(\mathrm{S}_{\alpha}\) 的  组成。）\(\mathbf{M}'\)\(y_{\alpha}'\)\(x + y_{\alpha}' \in \alpha\mathbf{M}\) 组成。）
\item
  设 A 为线性紧赋值环；令 K 为 A 的分式域。证明：若 M 是无挠 A-模，且 \(\mathbf{M}_{(\mathbf{K})}\) 有可数基，则 M 是一族可数个秩 1 的 A-模的直和。（将 M 看作递增的纯子模序列 \((\mathbf{M}_{i}^{\prime})\) 的并，其中 \(M_{i}^{\prime}\) 的秩为 i，并对每个 i 将 (b) 应用于 \(M_{i}^{\prime}\) 及其子模 \(M_{i-1}^{\prime}\)。）
\item
  在 (c) 中关于 A 与 M 的假设下，令 N 为 M 的有限秩纯子模。证明 N 是 M 的直因子（注意，每个作为有限个秩 1 的 A-模之直和的 A-模都是线性紧的，并使用 (b)）。
\item
  在 (b) 中关于 M 与  的相同假设下，证明对所有 \(M'\)，存在 \(x \in M\)，使得 \(y' \in M'\) 的湮没子等于元素 \(x' = x + y'\) 的湮没子。（令 a 为 \(x + M' \in M/M'\) 的湮没子；对所有 \(x + M'\)，令 \(\alpha \in a\) 为 M' 的子集，它由满足 \(T_{\alpha}\) 的  组成；证明 \(M'\)\(y'\) 的交非空。）\(\alpha(x + y') = 0\)\(T_{\alpha}\) 的交非空。）
\item
  设 A 为赋值环，并满足对每个理想 ，A-模 A/a（带离散拓扑）线性紧。证明每个有限生成的无挠 A-模 M 都是有限个单生成 A-模的直和。（若 \(a \neq 0\) 是 M 的一组生成元，则对 n 作归纳，取一个湮没子最小的 \((z_{i})_{1 \leqslant i \leqslant n}\)，并注意它生成的 M 的子模是纯的。）\(z_{i}\)，并注意它生成的 M 的子模是纯的。）
\end{enumerate}

\begin{enumerate}
\def\labelenumi{\arabic{enumi}.}
\setcounter{enumi}{6}
\tightlist
\item
  设 K 为域，v 为 K 上的离散赋值，且 K 关于 \(T_{v}\) 完备；令 A 与 p 分别表示 v 的环与理想，U = A - p 表示 A 的可逆元所成的集合。令 u 为 K 到 K 的一个子域上的同构。
\end{enumerate}

\begin{enumerate}
\def\labelenumi{(\alph{enumi})}
\item
  证明  且 \(u(\mathfrak{p}) \subset \mathbf{A}\)。（为证明第一点，注意对所有 \(u(\mathbf{U}) \subset \mathbf{U}\)，当 n > 0 与 v 的剩余域的特征互素时，方程 \(z \in \mathfrak{p}\) 由 Hensel 引理在 \(x^n = 1 + z\) 中有解；若 \(\mathbf{A}\)，则推出整数 \(n > 0\)\(v\)\(v(u(z)) < 0\) 将被每个与 v 的剩余域的特征互素的整数 \(v(u(z))\) 整除。然后由第一点推出第二个断言。）\(n > 0\)\(v\) 整除。然后由第一点推出第二个断言。）
\item
  从 (a) 推出：或者 ，或者 u 连续（考虑 v 的一个统一化元在 u 下的像）。\(u(\mathbf{K}^*) \subset \mathbf{U}\)\(u\)\(u\)\(v\)，或者 u 连续（考虑 v 的一个统一化元在 u 下的像）。
\item
  给出一个代数闭域 \(\Omega\) 的例子，使得取 \(\mathbf{K} = \Omega((\mathbf{X}))\) 时，\(\mathbf{K}\) 同构于 \(\mathbf{K}\) 的一个包含于 \(\mathbf{U} \cup \{0\}\) 的子域（参见代数，第 V 章 § 5，习题 13）。
\end{enumerate}

¶ 8. (a) 设 K 为域，v 为 K 上的高度 1 赋值，A 为其环。令 H 为 K 的紧子集（带拓扑 \(T_{v}\)），令 \(a \neq 0\) 为 K 的一点。证明存在无常数项的多项式 \(f \in K[X]\)，使得 \(f(a) = 1\)，\(f(H) \subset A\)。（证明 f 可取为如下形式的多项式

\[
1 - (1 - a ^ {- 1} \mathrm{X}) (1 - c _ {1} ^ {- 1} \mathrm{X}) ^ {n (1)} \dots (1 - c _ {r} ^ {- 1} \mathrm{X}) ^ {n (r)}
\]

其中  是  中适当选取的满足 \(c_{i}\)\(\mathbf{H}\) 的元素，而  是充分大的整数 \(v(c_{i}) < v(a)\)\(n(i)\)。 \(>0\)。 

\begin{enumerate}
\def\labelenumi{(\alph{enumi})}
\setcounter{enumi}{1}
\tightlist
\item
  令 X 为全不连通紧空间；赋予从 X 到 K 的连续映射环 \(\mathcal{C}(\mathbf{X};\mathbf{K})\) 以一致收敛拓扑。令 B 为 \(\mathcal{C}(\mathbf{X};\mathbf{K})\) 的包含常数且能分离 X 中各点的子环；证明 B 在 \(\mathcal{C}(\mathbf{X};\mathbf{K})\) 中稠密（使用 (a) 及一般拓扑，第 X 章 §4，习题 21 (b)）。
\end{enumerate}

¶ 9. 设 A 为离散赋值环，π 为 A 的统一化元，K 为 A 的分式域；假设 A 关于 π-进拓扑完备。内射 A-模（代数，第二章 §2，习题 11）与可除 A-模（代数，第七章 §2，习题 3）相同，并且都是同构于 K 或 K/A 的 A-模的直和（代数，第七章 §2，习题 3）。此外，每个单生成 A-模都同构于 K/A 的一个子模。A-模 Hom \(_{A}\) (M, K/A) 称为 A-模 M 的（代数）托里对偶，记作 M*；典范映射 c \(_{M}\) : M → M** 是单射（代数，第二章 §2，习题 13 (b)）。对于 M 的每个子模 N，M* 的由满足 u(x) = 0 的 u 组成的子模 N \(^{0}\) 称为 N 在 M* 中的正交补；M/N 的对偶典范地识别为 N \(^{0}\)，N 的对偶识别为 M*/N \(^{0}\)。直极限 lim M \(_{\alpha}\) 的托里对偶典范同构于逆极限 lim M \(_{\alpha}^{*}\)。

\begin{enumerate}
\def\labelenumi{(\alph{enumi})}
\tightlist
\item
  我们知道，秩 1 的 A-模（代数，第七章 § 4，
\end{enumerate}

习题 22）同构于以下形式之一的模：A、K、K/A 或 A/ \(\pi^{h}\) A。证明 K 与 A/ \(\pi^{h}\) A 的代数托里对偶分别同构于 K 与 A/ \(\pi^{h}\) A，A 的托里对偶同构于 K/A，而 K/A 的托里对偶同构于 A（利用 K 的 A-子模的知识以及 A 完备这一事实（相对于 K/A 的对偶））。推出：对每个有限秩 A-模 M，M* 是具有相同秩的模，且典范同态 \(c_{M}\) 为双射。

\begin{enumerate}
\def\labelenumi{(\alph{enumi})}
\setcounter{enumi}{1}
\item
  设 M 为 A-模，N 为 M 的有限秩子模。证明  中  的正交补通过 \(N^{0}\) 识别为 N（使用 (a)）。\(M^{**}\)\(c_{M}\) 识别为 N（使用 (a)）。
\item
  证明 Noether（分别 Artin）A-模 M 具有有限秩（将 M 嵌入其内射包中（代数，第二章 §2，习题 18））；于是 M* 是 Artin（分别 Noether）模。
\item
  令  表示 M* 上的拓扑，即 M 上的逐点收敛拓扑（赋予 K/A 离散拓扑）。证明：若 N 是 M 的子模，则 \(\sigma(M^{*}, M)\) 是由 \(M^{*}\)\(M\)\(K/A\)\(N\)\(M\)\(\sigma(N^{0}, M/N)\) 在  上诱导的拓扑，而 \(N^{0}\) 是拓扑  关于 \(\sigma(M^{*}, M)\) 的商拓扑。（关于第二点，注意若 P 是 M 的满足 \(\sigma(M^{*}/N^{0}, N)\)\(N^{0}\) 的子模，则 P/N 的对偶识别为 \(\sigma(M^{*}, M)\)\(P\)\(M\)\(N \subset P\)。）若 \(P/N\)\(N^{0}/P^{0}\)，则拓扑 \(M = \lim M_{\alpha}\) 是各拓扑 \(\sigma(M^{*}, M)\) 的逆极限。\(\sigma(M_{\alpha}^{*}, \overrightarrow{M_{\alpha}})\) 的逆极限。
\item
  拓扑  与 \(\sigma(\mathbf{K},\mathbf{K})\) 是 \(\sigma(\mathbf{A},\mathbf{K}/\mathbf{A})\)-进拓扑；拓扑 \(\pi\) 与 \(\sigma(\mathbf{A}/\pi^{h}\mathbf{A},\mathbf{A}/\pi^{h}\mathbf{A})\) 是离散拓扑。推出对每个 A-模 M，带有 \(\sigma(\mathbf{K}/\mathbf{A},\mathbf{A})\) 的模  是线性紧的（第三章 §2，习题 15；将 M 看作自由 A-模的商模）。\(\mathbf{M}^{*}\)\(\sigma(\mathbf{M}^{*},\mathbf{M})\) 是线性紧的（第三章 §2，习题 15；将 M 看作自由 A-模的商模）。
\item
  令 M、N 为两个 A-模；对于每个同态 ，令 \(u: \mathbf{M} \to \mathbf{N}\) 表示从 \(^{t}u\) 到 \(\operatorname{Hom}(u, \mathbf{l}_{\mathbf{K/A}})\)\(\mathbf{N}^{*}\) 的同态 \(\mathbf{M}^{*}\)，满足，满足
\end{enumerate}

\[
\left(^ {t} u (w)\right) (x) = w (u (x))
\]

对所有  与所有 \(x \in \mathbf{M}\)，证明 \(w \in \mathbf{N}^*\) 关于拓扑 \({}^t u\) 与 \(\sigma(\mathbf{N}^*, \mathbf{N})\) 连续。若 u 是 \(\sigma(\mathbf{M}^*, \mathbf{M})\) 的自同态 ，则 \(u\)\(x \mapsto \pi x\) 是 \(\mathbf{M}\)\({}^t u\) 的自同态 。对 \(w \mapsto \pi w\) 的每个子模 ，都有 \(\mathbf{M}^*\)\(\mathbf{P}\)。\(\mathbf{M}\)\((u(\mathbf{P}))^0 = {}^t u^{-1}(\mathbf{P}^0)\)。

¶ 10. 假设与记号同习题 9。

\begin{enumerate}
\def\labelenumi{(\alph{enumi})}
\item
  若 M 是离散拓扑 A-模，证明 M 是无挠模。若进一步 M 线性紧，证明 M 是 Artin 模（若 N 是自同态 \(x \mapsto \pi x\) 的核，注意 N 线性紧且离散，因而可看作向量 (A/ \(\pi\) A)-空间；然后使用第二章 § 2 的习题 20 (d)）。
\item
  由 (a) 推出每个线性紧拓扑 A-模都是严格线性紧的（第二章 § 2，习题 19）。
\item
  拓扑 A-模 M 的拓扑托里对偶是代数托里对偶  的子模 ，由从 M 到 K/A 的连续同态组成（后者赋予离散拓扑）。若 M 上的拓扑是线性的（第二章 § 2，习题 14），且 N 是 M 的闭子模，证明 M/N 的拓扑托里对偶识别为 \(M'\)，而 N 的拓扑托里对偶识别为 \(M^{*}\)\(N^{0} \cap M'\)（为确定 N 的对偶，注意对从 N 到 K/A 的每个连续同态 u，存在 M 的开子模 U，使得 u(x) = 0 对 \(M' / (N^{0} \cap M')\) 中的 x 成立，然后使用习题 9）。\(u(x) = 0\)\(N \cap U\) 中的 x 成立，然后使用习题 9）。
\item
  对于有限秩拓扑 A-模 M，拓扑托里对偶 \(M'\) 等于代数托里对偶 \(M^{*}\)（归结为秩 1 的模的情形）。
\item
  设 M 为拓扑线性的 Hausdorff 拓扑 A-模；证明对 M 中所有 ，存在 \(x \neq 0\) 使得 \(u \in M'\)（注意存在 M 的开子模 U 使得 \(u(x) \neq 0\)）；换言之，典范映射 \(x \notin U\) 是单射。推出若 N 是 M 的闭子模，则 \(\mathbf{M} \to (\mathbf{M}')^*\) 关于拓扑 \(N^0 \cap M'\) 在  中稠密（使用 (d)），从而在 \(N^0\) 中有 。\(\sigma(\mathbf{M}^*, \mathbf{M})\)\(N = M \cap (N^0 \cap M')^0\)\(M^{**}\)。
\end{enumerate}

证明 M 关于拓扑  在  中稠密。\((\mathbf{M}^{\prime})^{*}\)\(\sigma((\mathbf{M}^{\prime})^{*}, \mathbf{M}^{\prime})\) 中稠密。

\begin{enumerate}
\def\labelenumi{(\alph{enumi})}
\setcounter{enumi}{5}
\item
  设 M 为线性紧拓扑 A-模；证明典范嵌入 \(M \rightarrow (M')^{*}\) 是双射，且 M 上的拓扑（将 M 识别为 \((M')^{*}\) 后）等于 \(\sigma(M, M')\)。（注意若 U 是 M 的开子模，则由 (a) M/U 是 Artin 模，因而 \(U^{0}\) 是 Noether 模（习题 9 (c)），并推出 M 上所考虑的拓扑相同；借助 (e) 完成证明。）
\item
  设 M、N 为拓扑均为线性的两个拓扑 A-模；对于每个连续同态 ，有 \(u \colon \mathbf{M} \to \mathbf{N}\)，并且也（滥用记号地）用 \(^t u(\mathbf{N}') \subset \mathbf{M}'\) 表示从 \(^t u\) 到 \(\mathbf{N}'\) 的线性映射，其图与 \(\mathbf{M}'\) 相同。若 M 与 N 是 Hausdorff 的，证明 \(^t u\) 在 M 上的限制与 u 重合（分别将 M、N 看作典范嵌入 \(^t (^t u)\) 与 \(u\)\((\mathbf{M}')^*\) 中）；此外，对 N 的每个子模 Q，\((\mathbf{N}')\) 中）；此外，对 N 的每个子模 Q，
\end{enumerate}

\[
(\bar {u} ^ {1} (\mathbf {Q})) ^ {0} \cap \mathbf {M} ^ {\prime} = ^ {t} u (\mathbf {Q} ^ {0})
\]

（使用 (e)）。

\begin{enumerate}
\def\labelenumi{(\alph{enumi})}
\setcounter{enumi}{7}
\tightlist
\item
  设 M 为线性紧 A-模；由 (g) 和习题 9 (f) 推出 M 无挠的充要条件是 M' 可除，而 M 可除的充要条件是 M' 无挠。
\end{enumerate}

\exercisesection{6}

¶ 1. p-进域  的每个元素 z 都可唯一写成 \(Q_{p}\)，其中 \(p^{h}\sum_{k=0}^{\infty}c_{k}p^{k}\)、\(h\in Z\)，且对所有 \(c_{0}\neq0\) 有 （z 的“p-进展开”）。\(0\leqslant c_{k}<p\)\(k\geqslant0\)（z 的“p-进展开”）。

\begin{enumerate}
\def\labelenumi{(\alph{enumi})}
\tightlist
\item
  证明  的充要条件是存在两个整数 \(z \in \mathbf{Q}\)，使得对所有 \(m \geqslant 0, n \geqslant 1\) 都有 。（注意，若 \(c_{k+n} = c_k\)，其中 b 不被 p 整除，则\(k \geqslant m\)\(z = a / b \in \mathbf{Q}\)\(b\)\(p\)，其中 b 不被 p 整除，则
\end{enumerate}

\[
a - b \sum_ {k = 0} ^ {n} c _ {k} p ^ {k} = a _ {n + 1} p ^ {n + 1}
\]

其中 \(a_{n+1}\) 为整数，且序列 \(|a_k|\)（通常绝对值）有界。）

\begin{enumerate}
\def\labelenumi{(\alph{enumi})}
\setcounter{enumi}{1}
\tightlist
\item
  假设满足  的整数 k 的递增序列  满足 \((k_n)\)。证明 z 在 \(k\)\(c_k \neq 0\)\(\limsup_{n \to \infty} (k_{n+1}/k_n) = +\infty\) 上超越。（对 \(z\)\(\mathbf{Q}\)，令 \(x \in \mathbf{Q}\) 与 \(|x|\) 分别表示通常绝对值与 p-进绝对值。假设 \(|x|_p\) 是满足 \(p\)\(\mathrm{P} \in \mathbf{Z}[\mathrm{X}]\) 的不可约多项式；若 \(\mathrm{P}(z) = 0\)，使用 Taylor 公式证明 \(z_n = p^n \sum_{k=0}^{\infty} c_k p^k\) 随 n 趋于 \(|\mathrm{P}(z_n)/(z - z_n)|_p\) 时趋于一个  的极限；然后考虑通常绝对值 \(\neq 0\)，从假设得到矛盾。）\(n\)\(+\infty\)\(|\mathrm{P}(z_n)|\)，从假设得到矛盾。）
\end{enumerate}

\begin{enumerate}
\def\labelenumi{\arabic{enumi}.}
\setcounter{enumi}{1}
\tightlist
\item
  设 D 为特征 \(\neq2\) 的非交换域，Z 为其中心。假设 D 中每个包含 Z 的交换子域 K 在 Z 上的次数均 \(\leqslant2\)。
\end{enumerate}

\begin{enumerate}
\def\labelenumi{(\alph{enumi})}
\item
  不使用引理 1，证明 。（取 \([\mathbf{D}:\mathbf{Z}] = 4\) 使得 \(a\in \mathbf{D} - \mathbf{Z}\)，并对所有 \(a^2\in \mathbf{Z}\) 令 ；注意 \(\sigma (x) = axa^{-1}\) 分解为 Z 上两个向量子空间 \(x\in \mathbf{D}\)\(\mathbf{D}\) 与 \(\mathbf{Z}\)\(\mathrm{D}_{+}\) 的直和，使得 \(\mathrm{D}_{-}\) 在 \(\sigma (x) = x\) 上成立，而 \(\mathrm{D}_{+}\) 在 \(\sigma (x) = -x\) 上成立；还注意 \(\mathrm{D}_{-}\) 是 \(\mathrm{D}_{+}\) 的子域，而 \(\mathbf{D}\) 是 \(\mathrm{D}_{-}\) 上的 1 维向量空间；最后证明 \(\mathrm{D}_{+}\) 不可能不同于 \(\mathrm{D}_{+}\)。）\(\mathbf{Z}(a)\)。）
\item
  证明  是 Z 上的四元数代数。（借助 (a) 构造 \(\mathbf{D}\) 在 Z 上的一组基。）\(\mathbf{Z}\)\(\mathbf{D}\)\(\mathbf{Z}\) 在 Z 上的一组基。）
\end{enumerate}

\exercisesection{7}

\begin{enumerate}
\def\labelenumi{\arabic{enumi}.}
\item
  设 \((\phi_{t})_{t\in I}\) 为域 \(K\) 的一族位，其值取于有限个域；进一步假设对某个 \(x\in K\)，\(\phi_t(x)\) 所成集合是有限的。证明存在形如第 1 节 (1) 的多项式 \(f(X)\)，使得 \(f(x)\neq 0\)，且元素 \(z = f(x)^{-1}\) 具有如下性质：(i) \(\phi_t(x) = \infty\) 蕴含 \(\phi_t(xz) = 0\) 且 \(\phi_t(z) = 0\)；(ii) \(\phi_t(x)\neq \infty\) 蕴含 \(\phi_t(xz)\neq \infty\) 且 \(\phi_t(z)\neq 0\)。（证明与引理 1 相同。）
\item
  用第 1 节命题 1 的假设与记号，令 q 为 B 的素理想；证明 \(B_{q}\) 是赋值环。
\item
  令 （\(\mathbf{A}_i\)）为域 \(1 \leqslant i \leqslant n\) 的两两独立赋值环，并令 \(\mathbf{K}\)。对每个 \(\mathbf{A} = \bigcap_{i} \mathbf{A}_i\)，令 \(i\) 为 \(a_i\) 的理想 ，并记 \(\neq 0\)；\(\mathbf{A}_i\)\(a = \bigcap_{i} a_i\)；
\end{enumerate}

证明对所有  都有 \(i\)。反之，若 \(\mathfrak{a}_i = \mathbf{A}_i\mathfrak{a}\) 是 \(\mathfrak{b}\) 的理想 ，且 \(\neq 0\)，则 \(\mathbf{A}\)\(\mathfrak{b}_i = \mathbf{A}_i\mathfrak{b}\)（使用第 2 节的定理 1）。\(\mathfrak{b} = \bigcap_i \mathfrak{b}_i\)（使用第 2 节的定理 1）。

\exercisesection{8}

\begin{enumerate}
\def\labelenumi{\arabic{enumi}.}
\item
  设 K 为域，A = K{[}{[}X, Y, Z{]}{]} 为 K 上三个不定元的形式幂级数环；令 \(v'\)（分别 \(v''\)）为 A 上取值于按字典序排序的群 \(Z \times Z\) 的赋值，满足 \(v'(X) = (1, 0)\)、\(v'(Y) = (0, 1)\)，且 \(v'\) 在 K{[}{[}Z{]}{]} 上是非真赋值（分别地，\(v''(X) = (1, 0)\)、\(v''\) 在 K{[}{[}Y{]}{]} 上是非真赋值、\(v''(Z) = (0, 1)\)）。令 \(\sigma\) 为保持 K 与 X 不变且满足 \(\sigma(Y) = Z\)、\(\sigma(Z) = Y\) 的 A 的自同构；若 B 是 A 中由 \(\sigma\) 不变元素组成的子环，E（分别 F）为 A（分别 B）的分式域，则 \(v'\) 与 \(v''\)（典范延拓到 E 后）在 F 上的限制相同，F 关于 v 定义的拓扑完备，而 E 是 F 的二次扩张；E 上的两个赋值 \(v'\)、\(v''\) 相依但不等价。
\item
  设  为向 2-进域 \(K_{0}\) 添加所有多项式 \(Q_{2}\) 的根所得的域；令 v 为 K＿0 上延拓 2-进赋值的唯一赋值，并令 K 为 K＿0 关于 v 的完备化。证明多项式 \(X^{2^{n}} - 2\) 在 K\(K_{0}\)\(K_{0}\)\(X^{2} - 3\){[}X{]} 中不可约；令  为该多项式的根域，令 \(K'\) 为 v 到 \(v'\) 的延拓；则\(K'\) 的延拓；则
\end{enumerate}

\[
n = \left[ \mathrm{K} ^ {\prime}: \mathrm{K} \right] = 2, \quad e (v ^ {\prime} / v) = f (v ^ {\prime} / v) = 1.
\]

\begin{enumerate}
\def\labelenumi{\arabic{enumi}.}
\setcounter{enumi}{2}
\tightlist
\item
  令 k 为素域  上无穷多个不定元的有理函数域 ，并令 \(\mathbf{F}_{p}(\mathbf{X}_{n})_{n\in\mathbf{N}}\) 为 k 上两个不定元的有理函数域。\(F_{p}\)\(K = k(\mathbf{U}, \mathbf{V})\) 为 k 上两个不定元的有理函数域。
\end{enumerate}

\begin{enumerate}
\def\labelenumi{(\alph{enumi})}
\item
  证明形式幂级数域  的元素  在有理函数域 \(\mathbf{P}(\mathbf{U}) = \sum_{n=0}^{\infty} \mathbf{X}_n^p \mathbf{U}^{np}\) 上不代数。映射 \(k((\mathbf{U}))\)\(k(\mathbf{U})\) 从 \(\mathbf{F}(\mathbf{U}, \mathbf{V}) \mapsto \mathbf{F}(\mathbf{U}, \mathbf{P}(\mathbf{U}))\) 到 \(k[\mathbf{U}, \mathbf{V}]\) 可延拓为 \(k((\mathbf{U}))\) 到 \(\mathbf{K}\) 的一个子域上的同构；将 \(k((\mathbf{U}))\) 上等于形式幂级数的阶的赋值（§ 3，第 3 节，例 3）限制到该子域，得到 \(k((\mathbf{U}))\) 上的离散赋值 v，其剩余域为 k。\(v\)\(\mathbf{K}\)\(k\) 上的离散赋值 v，其剩余域为 k。
\item
  令  为  的代数扩张，使得 \(\mathbf{K}'\)\(\mathbf{K}(\mathbf{V}^{1/p})\)\(\mathbf{K}\)；若 \([\mathbf{K}':\mathbf{K}] = p\) 为 \(v'\) 上延拓 v 的唯一赋值，证明 \(\mathbf{K}'\)。因此，赋值 \(v\)\(e(v'/v) = f(v'/v) = 1\) 的环不是赋值 v 的环上的有限生成模。\(v'\)\(v\) 的环不是赋值 v 的环上的有限生成模。
\end{enumerate}

\begin{enumerate}
\def\labelenumi{\arabic{enumi}.}
\setcounter{enumi}{3}
\item
  设  为域，\(k\) 为 k 上两个不定元的有理函数域，v 为 \(\mathbf{K} = k(\mathbf{X}, \mathbf{Y})\) 上取值于按字典序排序的群 \(k\)\(v\)\(\mathbf{K}\) 的赋值，满足 \(\mathbf{Z} \times \mathbf{Z}\)、\(v(\mathbf{X}) = (0, 1)\)。令 \(v(\mathbf{Y}) = (1, 0)\) 为域 \(\mathbf{K}'\)；证明 v 到 \(\mathbf{K}(\sqrt{\mathbf{X}})\) 有唯一延拓 ，且 \(v\)\(v'\)、\(\mathbf{K}'\)\(e(v'/v) = 2\)，但赋值 \(f(v^{\prime}/v)=1\) 的环不是赋值 v 的环上的有限生成模。\(v^{\prime}\)\(v\) 的环不是赋值 v 的环上的有限生成模。
\item
  设 K 为域，A 为 K 的赋值环，L 为 K 的有限代数扩张。令  为 K 的另一个赋值环；令 \(A' \supset A\)（）为 L 中满足  的赋值环，并令 \(A_i' (1 \leqslant i \leqslant m)\)\(A_i' \cap K = A'\) 为它们各自的剩余域。令 k 为 \(k_i'\) 的剩余域，\(A'\) 为 A 的典范像所在的 K 的赋值环；最后，令 \(A''\)（）为  的满足 \(A_{ij}' (1 \leqslant j \leqslant n_i)\)\(k_i'\) 的赋值环，并令 \(A_{ij}'' \cap k = A''\) 为 \(A_{ij}\) 中 \(A_i'\) 的逆像（）。\(A_{ij}'' (1 \leqslant i \leqslant m, 1 \leqslant j \leqslant n_i)\)）。
\end{enumerate}

\begin{enumerate}
\def\labelenumi{(\alph{enumi})}
\item
  证明  是 \(\mathbf{A}_{ij}\) 的彼此不同的赋值环，满足 \(\mathbf{L}\)，并且 \(\mathbf{A}_{ij} \cap \mathbf{K} = \mathbf{A}\) 的每个满足 \(\mathbf{B}\)\(\mathbf{L}\) 的赋值环 \(\mathbf{B} \cap \mathbf{K} = \mathbf{A}\) 都等于某个 。\(\mathbf{A}_{ij}\)。
\item
  证明
\end{enumerate}

\[
e \left(\mathrm{A} _ {i j} / \mathrm{A}\right) = e \left(\mathrm{A} _ {i} ^ {\prime} / \mathrm{A} ^ {\prime}\right) e \left(\mathrm{A} _ {i j} ^ {\prime \prime} / \mathrm{A} ^ {\prime \prime}\right) \quad \text { and } \quad f \left(\mathrm{A} _ {i j} / \mathrm{A}\right) = f \left(\mathrm{A} _ {i j} ^ {\prime \prime} / \mathrm{A} ^ {\prime \prime}\right).
\]

¶ 6. 设 K 为域，v 为 K 上的赋值，A 为 v 的环。

\begin{enumerate}
\def\labelenumi{(\alph{enumi})}
\item
  假设 A 是 Hensel 的（第三章 § 4，习题 3）；此时我们也约定称 K 关于 v 是 Hensel 的。于是对 K 的每个代数扩张 L 以及 L 上延拓 v 的每个赋值 ，\(v'\) 的环  都是 Hensel 的（第三章 § 4，习题 4）。\(A'\)\(v'\) 都是 Hensel 的（第三章 § 4，习题 4）。
\item
  若 K 关于 v 完备且 v 的高度为 1，则 A 是 Hensel 的。给出 K 关于 v 完备、v 的高度为 2 而 A 不是 Hensel 环的例子（参见习题 1）。
\item
  若 K 关于 v 线性紧，证明 K 是 Hensel 的。（在第三章 § 4，习题 3 的条件 (H) 的记号中，令 L 表示 A 的素理想所成的集合（按包含关系）全序，令 \(\mathcal{L}\) 表示有序对 (B, φ) 所成的集合，其中 B 是 L 的良序子集（关于关系 ⊃），而 φ: p ↦ (Q\_p, Q'\_p) 是从 B 到集合 A{[}X{]} × A{[}X{]} 的映射，满足以下性质：(1) Q\_p 与 Q'\_p 分别为次数 deg( \(\overline{Q}\) ) 与 deg( \(\overline{Q}'\) ) 的首一多项式；(2) 若 p ⊃ q 是 B 的两个元素，则多项式 Q\_p - Q\_q 与 Q'\_p - Q'\_q 的系数属于 p；(3) 多项式 P - Q\_pQ'\_p 的系数属于 p；(4) \(\bar{f}(Q_{p}) = \overline{Q}, \bar{f}(Q'_{p}) = \overline{Q}'\)。在 L 上规定序：当 B ⊂ B' 且 φ' 是 φ 到 B' 的延拓时，置 (B, φ) ≤ (B', φ')。证明 L 有极大元 (B₀, φ₀)，且 B₀ 有一个等于 (0) 的末元：反设之，并根据 B₀ 是否有末元分两种情形。前一种情形中，若 p 是该末元，取 c ∈ A，使得 v(c) 是 v 在多项式 P - Q\_pQ'\_p 的系数集上所取的最小值，并令 p' 为包含 c 的 A 的最小素理想；若 p' = p，则使用 § 4 的习题 2 (a) 以及线性紧模被闭子模取商后仍线性紧、因而完备这一事实。反之，若 B₀ 没有末元，则直接使用 A 线性紧的假设。） 
\item
  用习题 5 的记号，证明 A 为 Hensel 环的充要条件是  与 \(\mathbf{A}'\) 都是 Hensel 环。（注意，若 \(\mathbf{A}''\)，则存在 \(\mathbf{P} \in \mathbf{A}'[\mathbf{X}]\)，使得 \(a \in \mathbf{A}\)（其中 n 是 \(a^n\mathbf{P}(\mathbf{X}/a) \in \mathbf{A}[\mathbf{X}]\) 的次数）；为验证 \(n\)\(\mathbf{P}\) 的第三章 § 4，习题 3 的公理 (H)，只需考虑 \(\mathbf{A}'\) 中的多项式；使用 \(\mathbf{A}[\mathbf{X}]\) 这一事实，其中 \(\mathbf{A}' = \mathbf{A}_{\mathfrak{p}}\) 是包含于极大理想 \(\mathfrak{p}\) 的  的素理想，并注意若 \(\mathbf{A}\) 中的两个多项式强互素，则它们在 \(\mathfrak{m}\)\((\mathbf{A}/\mathfrak{p})[\mathbf{X}]\) 中的像也强互素。）\((\mathbf{A}/\mathfrak{m})[\mathbf{X}]\) 中的像也强互素。）
\item
  假设 A 是 Hensel 环；于是对 K 的每个代数扩张 L，在 L 上存在（除等价外）唯一一个延拓 v 的赋值。（若  是首一不可约多项式，\(P \in A[X]\)（\(x_{i}\)）是 P 在其分裂域 N 中的不同根，证明对 N 上每个延拓 v 的赋值 \(1 \leqslant i \leqslant m\)，当 \(v'\) 时各  都相等。）反之，若赋值环 A 具有此性质，则 A 是 Hensel 环（注意在 A 在 L 中的整闭包 \(v'(x_{i})\) 中，位于 m 之上的极大理想至多有一个（第五章 § 1，习题 13），从而若 \(1 \leqslant i \leqslant m\)\(A'\) 是首一不可约多项式，则其在 \(P \in A[X]\) 中的像不可能是两个互素多项式的乘积）。\((\mathbf{A}/\mathfrak{m})[\mathbf{X}]\) 中的像不可能是两个互素多项式的乘积）。
\end{enumerate}

\begin{enumerate}
\def\labelenumi{\arabic{enumi}.}
\setcounter{enumi}{6}
\tightlist
\item
  设 K 为域，v 为 K 上的赋值，A 为 v 的环，m 为 v 的理想。假设 A 是 Hensel 环；令 L 为 K 的 n 次有限扩张；令 \(v'\) 为 L 上延拓 v 的唯一赋值（习题 6 (e)），A' 为其环，\(m'\) 为其理想。令 x 为 A' 的任意元素；证明 x 在 A'/m' 中的类 \(\bar{x}\) 在 A/m 上的次数整除 n，且 \(v'(x)\) 在 \(\Gamma_{v'}/\Gamma_v\) 中的类的阶也整除 n（考虑 x 在 K 上的极小多项式）。
\end{enumerate}

¶ 8. 设 K' 为域，B 为 K' 上赋值 v 的环，G 为 K' 的有限自同构群，K 为 K' 中由 G 不变元素组成的子域，且 A = K ∩ B 是 K 的赋值环，对应于赋值 w = v\textbar K。

\begin{enumerate}
\def\labelenumi{(\alph{enumi})}
\item
   上延拓 \(\mathbf{K}'\) 的赋值是 \(w\)（其中 \(v \circ \sigma\)）（第 6 节，命题 6 的推论 1），而 \(\sigma \in \mathcal{G}\) 在 \(\mathbf{A}'\)\(\mathbf{A}\) 中的整闭包 \(\mathbf{K}'\) 是所有  的交，其中 \(\sigma \cdot \mathbf{B}\) 遍历 \(\sigma\)。若 \(\mathcal{G}\) 是 \(\mathfrak{p}(\mathbf{B})\) 与极大理想 \(\mathbf{A}'\) 的交，则 \(\mathfrak{m}(\mathbf{B})\)。证明分解群 \(\sigma \cdot \mathfrak{p}(\mathbf{B}) = \mathfrak{p}(\sigma \cdot \mathbf{B})\)（第五章 § 2，第 2 节）是由满足 \(\mathcal{G}^{\mathbb{Z}}(\mathfrak{p}(\mathbf{B}))\) 的  组成的  的子群；记 \(\mathcal{G}\)，并令 \(\sigma\)\(\sigma \cdot \mathbf{B} = \mathbf{B}\)\(\mathcal{G}^{\mathbb{Z}} = \mathcal{G}^{\mathbb{Z}}(\mathfrak{p}(\mathbf{B}))\) 为 \(\mathbf{K}^{\mathbb{Z}}\) 中由 \(\mathbf{K}'\) 不变元素组成的子域，令 \(\mathcal{G}^{\mathbb{Z}}\) 为 v 在 \(\mathbf{B}^{\mathbb{Z}}\) 上诱导的赋值的环，它等于 \(\mathbf{K}^{\mathbb{Z}}\)；回忆 \(v\)\(\mathbf{B} \cap \mathbf{K}^{\mathbb{Z}}\) 与 \(\mathbf{B}^{\mathbb{Z}}\) 的剩余域相同，且极大理想 \(\mathbf{A}\) 由 \(\mathfrak{m}(\mathbf{B}^{\mathbb{Z}})\) 生成（第五章 § 2，第 2 节，命题 4）。\(\mathfrak{m}(\mathbf{A})\mathbf{B}^{\mathbb{Z}}\) 生成（第五章 § 2，第 2 节，命题 4）。
\item
  令  表示 v 到 \(v^{\mathbf{Z}}\) 的限制，令 \(v\)\(\mathbf{K}^{\mathbf{Z}}\) 与 \(\Gamma\) 表示 w 与 \(\Gamma^{\mathbf{Z}}\) 的阶群。证明 \(w\)\(v^{\mathbf{Z}}\)。（令 \(\Gamma = \Gamma^{\mathbf{Z}}\) 为 \(\mathcal{G}^{\mathrm{D}} \supset \mathcal{G}^{\mathbf{Z}}\) 的子群，它由满足 \(\mathcal{G}\) . B 依赖于 B（§ 7，第 2 节）的  组成。对 \(\sigma \in \mathcal{G}\) 的阶数作归纳。若 \(\sigma\)\(\mathcal{G}\)，令 \((\mathcal{G}: \mathcal{G}^{\mathrm{D}}) > 1\) 为 \(\mathbf{K}''\) 的不变域，令 \(\mathcal{G}^{\mathrm{D}}\)；证明 \(\mathbf{A}'' = \mathbf{K}'' \cap \mathbf{B}\) 的阶群等于 \(v|\mathbf{K}''\)：为此使用逼近定理（§ 7，第 2 节，定理 1），证明对所有 \(\Gamma\)，存在 \(x \in K''\)，使得 \(y \in K''\)，且相对于 K 的 y 的共轭元中除 y 外的各个 \(v(y) = v(x)\) 满足 ；于是可将 G 换为 \(v(y_i) = 0\)，并应用归纳假设。若 \(y_i\)\(y\)\(K\)\(y\)\(G\)\(G^D\)，令 \(G^D = G\) 为由 \(B_0\)（其中 \(K'\)\(\sigma.B\)）生成的 \(\sigma \in G\) 的赋值环；令  为其剩余域，\(\bar{K}' = \kappa(B_0)\) 为典范同态，\(\pi: B_0 \to \bar{K}'\) 为 \(\bar{B} = \pi(B)\) 的赋值环，\(\bar{K}'\) 为相应的 \(\bar{v}\) 上的赋值，使得在 \(\bar{K}'\) 中 ，从而 \(v = \bar{v} \circ \pi\)（\(B_0\)\(\bar{\Delta}\) 的阶群）是 \(\bar{v}\)（v 的阶群）的子群；若 \(\Delta\)，\(v\)\(\Delta_0 = \Delta/\bar{\Delta}\) 为典范同态，则 \(\bar{\omega}: \Delta \to \Delta_0\) 是对应于 \(v_0 = \bar{\omega} \circ v\) 且在 G 下不变的  上的赋值。取商后，G 在 \(K'\) 上定义一个自同构群 ；令 N 为典范同态 \(B_0\)\(G\)\(G\)\(\bar{G}\) 的核；证明 \(\bar{K}'\)\(N\)\(G \to \bar{G}\)，并推出若 \(N \subset G^Z\)，可将 G 换为 \(N \neq \{e\}\)，然后应用归纳假设。最后假设 \(G\)\(G/N\)；令 \(N = \{e\}\)，令 \(A_0 = K \cap B_0\) 为 \(w_0\) 到 K 的限制；证明 \(v_0\) 的阶群为 \(K\)\(w_0\)，且 \(\Delta_0\) 是 \(\pi(A_0) = \bar{K}\) 的不变域（参见第 1 节，引理 2）；若 \(\bar{G}\) 为 \(\bar{w}\) 到  的限制，\(\bar{K}\) 为其阶群，则 \(\bar{v}\)\(\bar{\Gamma}\) 典范同构于 \(\Delta_0\)。另一方面，群 \(\Gamma/\bar{\Gamma}\) 是 \(G^Z(p(\bar{B}))\) 在 \(G^Z(p(B))\) 中的典范像；注意 \(G^Z\)，因而可应用归纳假设证明 \(G^D \neq G^Z\) 等于 \(\bar{\Gamma}\) 到 \(F^Z(θ)\)\(\bar{v}\) 的限制的阶群 \(K^Z\) 的限制的阶群 
\item
  令  为  的惯性群（第五章 § 2，第 2 节），并令  表示 \(\mathcal{G}^{\mathrm{T}} = \mathcal{G}^{\mathrm{T}}(\mathfrak{p}(\mathrm{B}))\)\(\mathfrak{p}(\mathrm{B})\)\(\mathbf{K}^{\mathrm{T}}\) 中由 \(\mathbf{K}'\) 的不变元素组成的子域，令 \(\mathcal{G}^{\mathrm{T}}\) 表示 v 在 \(\mathbf{B}^{\mathrm{T}}\) 上诱导的赋值的环，它等于 \(\mathbf{K}^{\mathrm{T}}\)；回忆：若 \(v\)\(\mathbf{B} \cap \mathbf{K}^{\mathrm{T}}\) 与 \(k\) 分别为 A 与 B 的剩余域，则 \(k'\) 的剩余域  是 k 的拟 Galois 扩张 \(k^{\mathrm{T}}\) 中的最大可分扩张，并且 \(\mathbf{B}^{\mathrm{T}}\)\(k\)\(k'\) 典范同构于 \(k\)\(\mathcal{G}^{\mathrm{Z}} / \mathcal{G}^{\mathrm{T}}\)（或 \(k\)\(k'\)）的 \(k^{\mathrm{T}}\) -自同构群（第五章 § 2，定理 2 与命题 5）。令  为 v 到 \(v^{\mathrm{T}}\) 的限制；证明 \(v\)\(\mathbf{K}^{\mathrm{T}}\) 的阶群也等于 \(v^{\mathrm{T}}\)（将引理 2（第 1 节）应用于 \(\Gamma\) 对 \(\mathbf{K}^{\mathrm{T}}\) 的扩张）。\(\mathbf{K}^{\mathrm{Z}}\) 的扩张）。
\end{enumerate}

\begin{enumerate}
\def\labelenumi{\arabic{enumi}.}
\setcounter{enumi}{8}
\tightlist
\item
  设 K 为域，L 为 K 的有限代数扩张，\(v'\) 为 L 上的赋值，A' 为其环，v 为 \(v'\) 到 K 的限制，\(A = A' \cap K\) 为其环。
\end{enumerate}

\begin{enumerate}
\def\labelenumi{(\alph{enumi})}
\item
  假设 A 是 Hensel 环（习题 6）；证明乘积  整除 \(e(v'/v)f(v'/v)\)，且商 \(n = [\mathbf{L}:\mathbf{K}]\) 是 v 的剩余域 k 的特征指数的幂。（将其归结为 L 是 K 的 Galois 扩张的情形；使用习题 6 (e)、习题 7 以及第五章 §2 的定理 2 与命题 5。）\(n / e(v'/v)f(v'/v)\)\(k\)\(v\) 是 v 的剩余域 k 的特征指数的幂。（将其归结为 L 是 K 的 Galois 扩张的情形；使用习题 6 (e)、习题 7 以及第五章 §2 的定理 2 与命题 5。）
\item
  进一步假设 n 不被 v 的剩余域 k 的特征指数整除，且 。证明：对于群 \(n\)\(k\)\(v\)\(f(v' / v) = 1\) 中元素  的阶所等于的每个整数 m，都存在 \(m\)\(\gamma\)，使得 \(\Gamma_{v'} / \Gamma_v\)\(x \in \mathbf{L}\) 的类等于 \(v'(x) \mod \Gamma_v\)，且 \(\gamma\)。\(x^m \in \mathbf{K}\)。
\end{enumerate}

¶ 10. 令 w 为 2-进域  上的 2-进赋值；在有理函数域 \(Q_{2}\) 上考虑延拓 w 且满足\(\mathbf{Q}_2(\mathbf{X})\)\(v\)\(w\) 上考虑延拓 w 且满足

\[
v (a _ {0} + a _ {1} \mathrm{X} + \dots + a _ {n} \mathrm{X} ^ {n}) = \inf _ {i} (w (a _ {i}))
\]

（§ 10，第 1 节，引理 1）；令 K 为  关于此赋值的完备化，并仍以 v 表示其赋值。令 L 为多项式 \(\mathbf{Q}_2(\mathbf{X})\) 在 K\(v\)\((\mathrm{Y}^2 - \mathrm{X})^2 - 2\){[}Y{]} 中的分裂域，令 \(v'\)\(v\) 为 L 上延拓 v 的唯一赋值。证明 {[}L: K{]} = 8，，\(e(v'/v) = 4\)；若 \(f(v'/v) = 2\)（分别 \(k\)）是 v（分别 \(k'\)）的剩余域，则 \(v\)\(v'\) 是 k 的二次根式扩张；证明不存在 L 的子扩张 E 满足 \(k'\)\(k\){[}E: K{]} = 2 且  到 E 的限制的剩余域等于 \(v'\)。（必然有 \(k'\)，其中 \(\mathrm{E} = \mathrm{K}(\sqrt{\alpha})\) 不是 K 中的平方，但满足 \(\alpha \in \mathrm{K}\)（模 2）；借助 L 在 \(\alpha \equiv \mathrm{X}\) 上的适当基表示 ，并注意 \(\sqrt{\alpha}\) 中不存在平方同余于 X 模 2 的元素。）\(\mathrm{K}(\sqrt{2})\)\(\mathrm{K}(\sqrt{2})\) 中不存在平方同余于 X 模 2 的元素。）

¶ 11. 设 K 为域，L 为 K 的有限 Galois 扩张，G 为其 Galois 群。令 v 为 L 上的赋值，其剩余域为 k、阶群为 Γ；假设 v 在 G 下不变，且 v\textbar K 的限制与 v 具有相同的剩余域 k，于是 \(G = G^{Z} = G^{T}\)（习题 8 的记号）。

\begin{enumerate}
\def\labelenumi{(\alph{enumi})}
\item
  令 、\(x \in \mathbf{L}^*\)。证明 \(\sigma \in \mathcal{G}\) 是 v 的单位，且它在 \(x^{-1}\sigma(x)\) 中的像  只依赖于赋值 \(v\)\(\varepsilon_{\sigma}(x)\) 以及 \(k^*\)\(v(x)\) 模 \(\sigma\) 的换位子群  的类，因而定义了一个 Z-双线性映射（也记作 \(\mathcal{G}'\)）\(\mathcal{G}\)\(\varepsilon\)，它在 \((\mathcal{G}/\mathcal{G}') \times \Gamma \to k^*\) 上等于 1。\((\mathcal{G}/\mathcal{G}') \times v(\mathbf{K}^*)\) 上等于 1。
\item
  对所有 ，令 \(\sigma \in \mathcal{G}\) 为从 \(\phi(\sigma)\) 到  的同态，它将  映到 \(\Gamma / v(\mathbf{K}^*) \to k^*\)\(\varepsilon_{\sigma}(x)\)（对所有 \(v(x)\)）；于是定义了同态\(x \in \mathbf{L}^*\)）；于是定义了同态
\end{enumerate}

\[
\phi \colon \mathcal {G} \to \operatorname{Hom} _ {\mathbf {Z}} (\Gamma / v (\mathrm{K} ^ {*}), k ^ {*}).
\]

证明：若 p 是 k 的特征指数，则  的核  的阶是 p 的幂。（否则， 中存在 \(p\)\(k\)\(\mathcal{N}\) 以及素数 \(\phi\)\(p\)\(\sigma \neq e\)，使得 \(\mathcal{N}\)\(q \neq p\)；对 L - K 中某个整元 x，写成 \(\sigma^q = e\)，其中 \(x\)\(L - K\)\(\sigma(x) = x + y\)，并计算 \(v(y) > v(x)\) 得到矛盾。）推出 \(\sigma^q(x)\) 是可解群。\(\mathcal{G}\) 是可解群。

\begin{enumerate}
\def\labelenumi{(\alph{enumi})}
\setcounter{enumi}{2}
\tightlist
\item
  假设  与 p 互素；证明 \(n = [\mathbf{L} : \mathbf{K}]\) 是双射，\(p\)\(\phi\)，并且若 \(e(\mathbf{L} / \mathbf{K}) = n\) 是 \(\nu\) 中元素的阶的最小公倍数，则 \(lcm\)\(\mathcal{G}\) 含有 \(k^*\) 次单位根。（使用 (b) 及第 1 节的引理 2。）\(\nu\) 次单位根。（使用 (b) 及第 1 节的引理 2。）
\end{enumerate}

¶ 12. 设 K 为域，v 为 K 上的赋值，A 为赋值 v 的环，Γ 为其阶群。

\begin{enumerate}
\def\labelenumi{(\alph{enumi})}
\tightlist
\item
  令 \(f(\mathbf{X}) = a_{0}\mathbf{X}^{n} + a_{1}\mathbf{X}^{n-1} + \cdots + a_{n}\) 为 A{[}X{]} 中满足 \(v(a_{0}) = 0\) 且所有根都属于 K 的多项式；令 \(x_{i} (1 \leqslant i \leqslant r)\) 为这些不同的根，令 \(k_{i}\) 为 \(x_{i} (1 \leqslant i \leqslant r)\) 的重数。令
\end{enumerate}

\[
g (\mathbf {X}) = b _ {0} \mathbf {X} ^ {n} + b _ {1} \mathbf {X} ^ {n - 1} + \dots + b _ {n} = b _ {0} (\mathbf {X} - y _ {1}) \dots (\mathbf {X} - y _ {n})
\]

另一个满足  且  属于 K 的 A{[}X{]} 中的多项式。证明存在 \(v(b_{0}) = 0\)\(y_{h}\)\(\lambda_{0} \in \Gamma\)，使得对每个不同指标 i、j 的有序对都有 \(v(x_{i} - x_{j}) < \lambda_{0}\)，并且对所有 \(\lambda \geqslant \lambda_{0}\)，存在 \(\mu \geqslant \lambda\)，使得对所有 i 满足 \(v(a_{i} - b_{i}) \geqslant \mu\) 时，对 \(1 \leqslant i \leqslant r\)，恰有 \(k_{i}\) 个指标 h 满足 \(v(x_{i} - y_{i}) \geqslant \lambda\)。（先用两种方法计算 \(v(f(y_{h}))\)，证明必有某个 i 使 \(v(x_{i} - y_{h}) \geqslant \lambda\)；然后计算

\[
v \left(\frac {f ^ {\prime} (z)}{f (z)} - \frac {g ^ {\prime} (z)}{g (z)}\right)
\]

在适当的点  处的值。）若进一步除 \(z \in \mathbf{K}\) 外均有 ，证明只要 \(b_{i} = a_{i}\) 充分大，\(i = n - 1\)\(\lambda_0\) 就两两不同（假设 \(y_{h}\) 不是 g 的单根，计算 ）。\(f'(y_h)/f(y_h)\)\(y_{h}\)\(g\)）。

\begin{enumerate}
\def\labelenumi{(\alph{enumi})}
\setcounter{enumi}{1}
\tightlist
\item
  以下假设 \(\mathbf{K}\) 是子域 \(\mathbf{K}_0\) 的代数闭包，且 \(\mathbf{A}_0 = \mathbf{A} \cap \mathbf{K}_0\) 是 Hensel 环。假设多项式 \(f\) 属于 \(\mathbf{A}_0[\mathbf{X}]\)，且在 \(\mathbf{K}_0\) 上不可约且可分。令 \(z \in \mathbf{K}\) 满足
\end{enumerate}

\[
v (z - v _ {i}) > v (z - x _ {j})
\]

对所有 \(j \neq i\) 成立；证明 \(\mathrm{K}_0(x_i) \subset \mathrm{K}_0(z)\)。（证明 \(z - x_i\) 关于 \(\mathrm{K}_0(x_i)\) 的所有共轭元都相等。）

\begin{enumerate}
\def\labelenumi{(\alph{enumi})}
\setcounter{enumi}{2}
\tightlist
\item
  假设  可分，并令 \(f \in \mathbf{A}_0[\mathbf{X}]\) 为 f 在 \(f = a_0 \prod_{i=1}^{\infty} f_i\) 中分解为首一不可约多项式的分解（事实上它们属于 \(f\)\(\mathbf{K}_0[\mathbf{X}]\)，参见第五章 § 1，第 3 节，命题 11）。证明（沿用 (a) 的记号）：若 \(\mathbf{A}_0[\mathbf{X}]\) 取充分大，则 g 分解为首一\(\mu\)\(g\) 取充分大，则 g 分解为首一
\end{enumerate}

可约因子在 \(K_{0}[X]\) 中可写成 \(b_{0}\prod_{i=1}^{n}g_{i}\)，其中对所有 i，\(g_{i}\) 与 \(f_{i}\) 次数相同（称为与 f 同类型的分解），并且对所有 i，\(f_{i}\) 与 \(g_{i}\) 的分裂域相同。（先证明 g 可分；再考虑一个包含 f 与 g 的根的有限 Galois 扩张 \(K_{0}\)，考察该扩张的 Galois 群如何置换这些根；最后使用 (b)。）

\begin{enumerate}
\def\labelenumi{(\alph{enumi})}
\setcounter{enumi}{3}
\tightlist
\item
  给出一个 f 不可约但不可分的例子，并且对所有 ，存在多项式\(f\)\(\mu \geqslant \lambda_0\)，存在多项式
\end{enumerate}

\[
g (\mathbf {X}) = b _ {0} \mathbf {X} ^ {n} + b _ {1} \mathbf {X} ^ {n - 1} + \dots + b _ {n} \in \mathrm{A} _ {0} [ \mathbf {X} ]
\]

使得对所有 i 有 \(v(a_{i} - b_{i}) \geqslant \mu\)，但它不可约（取 \(K_{0}\) 使得 \(K \neq K_{0}\)，且 K 是 \(K_{0}\) 的根式扩张）。

¶ 13. (a) 设 K 为域，v 为 K 上的非真赋值。令 E 为由超滤子 U 滤过的集合，\(\xi \mapsto x_{\xi}\)、\(\xi \mapsto y_{\xi}\) 为 E 到带拓扑 \(\mathcal{T}_{v}\) 的 K 的两个映射。证明：若映射 \(\xi \mapsto x_{\xi}y_{\xi}\) 关于 U 在 K 中收敛于 0，则两个映射 \(\xi \mapsto x_{\xi}\)、\(\xi \mapsto y_{\xi}\) 中至少有一个也如此。（注意，若 \(\xi \mapsto x_{\xi}\) 关于 \(\mathfrak{U}\) 没有聚点，则映射 \(z \mapsto x_{\xi}^{-1}\) 在 \(\mathfrak{U}\) 的某个子集上有界。）

\begin{enumerate}
\def\labelenumi{(\alph{enumi})}
\setcounter{enumi}{1}
\item
  令  为 \(f\) 中的非常数多项式。证明：若 \(\mathbf{K}[\mathbf{X}]\) 关于 \(\xi \mapsto f(x_{\xi})\) 在  中收敛于 0，则 \(\mathbf{K}\) 在 \(\mathfrak{U}\)\(\xi \mapsto x_{\xi}\) 中收敛（将 f 分解为 \(\hat{\mathbf{K}}\) 的代数扩张中的因子，并使用 (a)）。\(f\)\(\mathbf{K}\) 的代数扩张中的因子，并使用 (a)）。
\item
  假设 K 在  中代数闭。则 K 到自身的映射 \(\hat{\mathbf{K}}\) 是闭映射（使用 (b)）。推出：若 \(x \mapsto f(x)\) 是 \(g\) 中的有理函数，B 是 K 的有界闭子集，则 \(\mathbf{K}(\mathbf{X})\)（其中 P 是 g 在 B 中的极点集）在 K 中闭。\(g(\mathbf{B} - \mathbf{P})\)\(g\)（其中 P 是 g 在 B 中的极点集）在 K 中闭。
\item
  假设 \(\mathbf{K}\) 的代数闭包是 \(\mathbf{K}\) 的根式扩张。证明 \(\hat{\mathbf{K}}\) 是代数闭域。（将 (b) 应用于 \(\hat{\mathbf{K}}\) 的代数闭包及多项式 \(f \in \hat{\mathbf{K}}[\mathbf{X}]\)；还要使用习题 12 (a)。）
\end{enumerate}

¶ 14. (a) 域 K 关于赋值 v 为 Hensel 的充要条件是：K 是包含于  中的 K 的最大可分代数扩张，且 \(\hat{K}\) 是 Hensel 的。（为看出条件必要，使用习题 12 (b)；为证明充分，注意为验证第三章 § 4，习题 3 的条件 (H)，只需考虑 P 在 K 上可分的情形；因为若 Q 是 P 在 \(\hat{K}\) 中的不可约因子，且 \(\hat{K}[X]\) 属于 K\(Q^{p^{e}}\){[}X{]}，则 Q 是 \(Q^{p^{e}}\) 中 P 与  的最大公因子。）考虑 v 的高度为 1 的情形（参见习题 6 (b)）。\(\hat{K}[X]\) 的最大公因子。）考虑 v 的高度为 1 的情形（参见习题 6 (b)）。

\begin{enumerate}
\def\labelenumi{(\alph{enumi})}
\setcounter{enumi}{1}
\item
  由 (a) 推出 p-进域  中存在一个关于 p-进赋值为 Hensel 的可数子域。\(p\)\(\mathbf{Q}_p\)\(p\) 中存在一个关于 p-进赋值为 Hensel 的可数子域。
\item
  给出一个关于离散赋值为 Hensel 但不完备的域 K，使得 \(\hat{K}\) 是 K 的有限根式扩张（参见第五章 § 1，习题 20）。
\end{enumerate}

¶ 15. (a) 设 K 为域，\(v_1\)、\(v_2\) 为 K 上的两个独立赋值（§ 7，第 2 节），\(K_1\)、\(K_2\) 分别为 K 关于 \(v_1\)、\(v_2\) 的完备化。令 \(L_1\)、\(L_2\) 为满足 \(K \subset L_1 \subset K_1\)、\(K \subset L_2 \subset K_2\) 的两个域，它们是 Hensel 的（分别通过 \(v_1\)、\(v_2\) 的连续延拓）。令 \(g_1 \in L_1[X]\)、\(g_2 \in L_2[X]\) 为次数相同的两个可分多项式，次数为 n。证明存在多项式 \(h \in K[X]\)，使得 h 在 \(L_i[X]\) 中的分解类型（习题 12 (c)）与 \(g_i\) 相同（\(i = 1, 2\)）。（将其归结为 \(g_1\)、\(g_2\) 属于 K{[}X{]} 且为首一的情形；考虑多项式

\[
a ^ {n} g _ {1} (\mathbf {X} / a) + b ^ {n} g _ {2} (\mathbf {X} / a) - \mathbf {X} ^ {n},
\]

其中 \(v_{1}(a) \leqslant 0\)，\(v_{2}(a) > 0\) 任意大，且 \(v_{1}(b) > 0\) 任意大。）

\begin{enumerate}
\def\labelenumi{(\alph{enumi})}
\setcounter{enumi}{1}
\item
  由 (a) 推出：若 K 关于 \(v_{1}\) 为 Hensel，则 \(L_{2}\) 的代数闭包是 \(L_{2}\) 的根式扩张，且 \(K_{2}\) 是代数闭域（令 \(g_{2}\) 不可约，令 \(g_{1} \in K[X]\) 为 n 个不同的素次数因子的乘积）。
\item
  由 (b) 推出：若 \(\mathbf{K}\) 关于 \(v_{1}\) 与 \(v_{2}\) 都为 Hensel，则 \(\mathbf{K}\) 的代数闭包是 \(\mathbf{K}\) 的根式扩张。
\item
  证明若 K 关于离散赋值 v 为 Hensel，则 K 的代数闭包不可能是 K 的根式扩张，因而 K 不可能关于任何独立于 v 的赋值为 Hensel（使用代数，第五章 § 11，习题 12）。
\end{enumerate}

\begin{enumerate}
\def\labelenumi{\arabic{enumi}.}
\setcounter{enumi}{15}
\tightlist
\item
  设 K 为带有高度 1 赋值 v 的 Hensel 域。若 L 是 K 的无限可分代数扩张，证明 L 不可能关于延拓 v 的赋值（也记作 v）完备。（构造 L 的元素序列 ，使得 \((x_{p})\) 关于 K 的次数  趋于 \(n_{p}\)，且 \(x_{p}\)\(+\infty\) 严格大于 \(v(x_{p+1}-x_{p})\) 与其关于 K 的共轭元之间的差的赋值；证明序列 \(x_{p}\) 是不能在 E 中收敛的 Cauchy 序列，使用习题 12 (b)。）\((x_{p})\) 是不能在 E 中收敛的 Cauchy 序列，使用习题 12 (b)。）
\end{enumerate}

¶ 17. 设 K 为带有一个非代数闭赋值 v 的 Hensel 域，L 为 K 的满足 {[}K: L{]} 有限的子域。证明在这些条件下 L 关于 v 的限制是 Hensel 的。（反设不然，则存在 v|L 到 K 的代数闭包  的两个延拓 ，使得 \(v_{1}, v_{2}\)、\(v|L\)\(\Omega\)\(v_{1}\) 到 L 的某个有限代数扩张 L' 的限制不等价。推出在包含 K 的 L 的某个有限拟 Galois 扩张 K' 上存在两个不等价的赋值 \(v_{2}\)，且 K' 关于它们是 Hensel 的；使用习题 15 (c)，证明若 \(L'\)\(K'\)\(v_{1}', v_{2}'\)（p 是 \(K'\)\(E = L^{p^{-\infty}}\) 的特征指数），则 \(\Omega\)，但 \(E \neq \Omega\) 是 E 的有限扩张；借助代数，第六章 § 2，习题 31 完成论证。）\(\Omega\) 是 E 的有限扩张；借助代数，第六章 § 2，习题 31 完成论证。）

¶ 18. 设 K 为带有赋值 v 的 Hensel 域，k 为 v 的剩余域，并假设 k 的每个有限代数扩张都是 k 上的循环扩张（例如 k 有限时即如此）。令 A 为 v 的环，令 \(f \in A[X]\) 为首一多项式，并假设当 \(f(\mathbf{X}) = (\mathbf{X} - \alpha_{1}) \ldots (\mathbf{X} - \alpha_{n})\)（其中 \(\alpha_{i}\) 属于 K 的代数闭包）时，A 中的元素 \(D = \prod_{i < j} (\alpha_{i} - \alpha_{j})^{2}\)（f 的“判别式”）满足 \(v(D) = 0\)。令 \(\bar{f}_{j} (1 \leqslant j \leqslant s)\) 为 f 在 k\(\bar{f}\){[}X{]} 中的典范像  的不可约因子，令 \(r_{j}\) 为 \(\bar{f}_{j}\) 的次数。证明 f 的分裂域在 K 上的 Galois 群（看作 \(\alpha_{i}\) 的置换群）由一个置换 \(\sigma\) 生成，该置换分解为长度分别为 \(r_{1}, r_{2}, \ldots, r_{s}\) 的循环（注意 k 的给定次数扩张除 k-同构外唯一）。推出 D 是 A 中平方的充要条件是 n - s 为偶数（“Stickelberger 定理”；考察 D 在 \(\sigma\) 下不变的条件）。

\begin{enumerate}
\def\labelenumi{\arabic{enumi}.}
\setcounter{enumi}{18}
\tightlist
\item
  设 K 为带有赋值 v 的 Hensel 域。令 E 为 K 上的有限维向量空间，Q 为 E 上的非退化二次型，且关系 \(Q(x) = 0\) 蕴含 x = 0。令
\end{enumerate}

\[
\Phi (x, y) = \mathrm{Q} (x + y) - \mathrm{Q} (x) - \mathrm{Q} (y)
\]

为相伴的双线性型。证明 \(2v(\Phi(x,y)) \geqslant v(Q(x)) + v(Q(y))\)；推出

\[
v (\mathrm{Q} (x + y)) \geqslant \inf (v (\mathrm{Q} (x)), v (\mathrm{Q} (y))).
\]

¶ 20. 设 K 为特征 ≠2 的域，关于赋值 v 为 Hensel；令 ξ ↦ ξ̄ 为 K 的满足 v(ξ̄) = v(ξ) 的对合自同构，令 Φ 为 K 上有限维向量空间 E 上的非退化 Hermite 型。

\begin{enumerate}
\def\labelenumi{(\alph{enumi})}
\item
  令  为 \(U = (\alpha_{ij})\) 关于 E 的一组基 \(\Phi\) 的矩阵；证明存在 v 的阶群中的元素 \((e_i)\)，使得若 \(\lambda\) 是 E 上的另一个 Hermite 型，其关于 \(v\)\(\Phi'\) 的矩阵  对每个有序对 \(U' = (\alpha_{ij}')\) 都满足 ，则 \((e_i)\)\(v(\alpha_{ij} - \alpha_{ij}') > \lambda\) 与 \((i,j)\)\(\Phi\) 等价。（仿照代数第九章 § 6，习题 6，并使用上面的习题 14 (a)。）\(\Phi'\) 等价。（仿照代数第九章 § 6，习题 6，并使用上面的习题 14 (a)。）
\item
  假设  是恒等映射（因而 \(\xi \mapsto \overline{\xi}\) 是对称型），且赋值 \(\Phi\) 离散、赋范并满足 \(v\)；令 \(v(2) = 0\) 为 \(\pi\) 的统一化元。证明 \(v\) 除等价外由其指标 \(\Phi\) 以及分别定义在向量空间 \(\nu\)、\(\Psi_1, \Psi_2\)\(k^r\) 上的两个指标为 0 的对称双线性型 \(k^s\) 所刻画，其中  为 \(k\) 的剩余域，且 \(v\)。（借助 Witt 分解将问题归结为 \(r + s = n - 2\) 的指标为 0 的情形；借助习题 19，证明对所有 \(\Phi\)，由满足 \(i \geqslant 0\) 的  组成的集合  是 v 的环 A 上的模。证明若 \(\mathbf{M}_i\)、\(x \in \mathbf{E}\)\(v(\Phi(x, x)) \geqslant i\)\(v\)\(x \in \mathbf{M}_i\)，则 \(y \in \mathbf{M}_{i+1}\)，因而取商后可由 \(v(\Phi(x, y)) \geqslant i + 1\) 在向量 \(\Phi\)-空间 \(k\) 与 \(\mathbf{M}_0 / \mathbf{M}_1\) 上导出对称双线性型。使用 §3 的习题 6 (c) 以及方程 \(\mathbf{M}_1 / \mathbf{M}_2\) 在 \(\xi^2 = \alpha\) 时在 A 中有解这一事实。）考虑 k 有限的情形（参见代数第九章 §6，习题 4）。\(\alpha \equiv 1 (\text{mod. } \pi)\)\(k\) 时在 A 中有解这一事实。）考虑 k 有限的情形（参见代数第九章 §6，习题 4）。
\end{enumerate}

\begin{enumerate}
\def\labelenumi{\arabic{enumi}.}
\setcounter{enumi}{20}
\tightlist
\item
  设 K 为域，v 为 K 上的离散赋值，A 为 v 的环，\(\pi\) 为 v 的统一化元，k 为 v 的剩余域。
\end{enumerate}

\begin{enumerate}
\def\labelenumi{(\alph{enumi})}
\item
  令 P、R 为 A{[}X{]} 中的两个多项式，P 为首一，并假设：(1) ，其中 \(\deg(R) < h \cdot \deg(P)\) 为整数；(2) P 在 k\(h \geqslant 1\)\(\overline{P}\){[}X{]} 中的典范像  不可约。证明若多项式  在 \(Q = P^{h} + \pi R\) 中可约，则必有 h \(\widehat{K}[X]\)\textgreater{} 1，且 R 在 k\(\overline{P}\)\(\overline{R}\){[}X{]} 中的典范像  被  整除。推出：给定首一不可约多项式  及整数 \(r(X) \in k[X]\)，存在首一不可约可分多项式 \(h \geqslant 1\)，使得其典范像 \(Q \in A[X]\) 等于 \(\overline{Q}\)。\(r^{h}\)。
\item
  设  为 k 的次数 m 的代数扩张，h 为整数 \(k' = k(\alpha)\)。证明存在 K 的次数 hm 的代数扩张 L，使得 L 上除等价外只有一个延拓 v 的赋值 \(k\)\(m\)\(h\)\(\geqslant 1\)，且 \(L\)\(K\)\(hm\)\(v'\)、\(L\)\(v\)\(e(v'/v) = h\)，并且 \(f(v'/v) = m\) 的剩余域同构于 \(v'\)（使用 (a)）。\(k'\)（使用 (a)）。
\end{enumerate}

\begin{enumerate}
\def\labelenumi{\arabic{enumi}.}
\setcounter{enumi}{21}
\tightlist
\item
  \begin{enumerate}
  \def\labelenumii{(\alph{enumii})}
  \tightlist
  \item
    若域 \(\mathbf{K}\) 上存在离散赋值，证明 \(\mathbf{K}\) 的代数闭包在 \(\mathbf{K}\) 上是无限扩张。
  \end{enumerate}
\end{enumerate}

\begin{enumerate}
\def\labelenumi{(\alph{enumi})}
\setcounter{enumi}{1}
\tightlist
\item
  设 K 为域  的有限生成扩张。证明若 K 在 \(K_{0}\) 上不代数，则 K 上存在离散赋值 v，使得在 \(K_{0}\) 上 。\(v(x) = 0\)\(K_{0}\)。
\end{enumerate}

\exercisesection{9}

\begin{enumerate}
\def\labelenumi{\arabic{enumi}.}
\tightlist
\item
  \begin{enumerate}
  \def\labelenumii{(\alph{enumii})}
  \tightlist
  \item
    设 K 为域（不必交换），T 为 K 上与环结构相容的局部紧拓扑；证明 T 与 K 的域结构相容。（使用 R. Ellis 定理（一般拓扑，第 X 章 § 3，习题 25），或直接仿照积分论第七章 § 1，第 10 节及本段命题 1 的证明。）
  \end{enumerate}
\end{enumerate}

\begin{enumerate}
\def\labelenumi{(\alph{enumi})}
\setcounter{enumi}{1}
\tightlist
\item
  给出交换域 K 上的一个局部紧拓扑的例子，它与 K 的加法群结构相容，但不与其环结构相容。（取 K 为紧整环 A 的分式域（例如 A 为形式幂级数环 \(k[[X, Y]]\)，其中 k 有限），并取 K 中 0 的邻域在 A 中的 0 的邻域作为基本邻域系。）
\end{enumerate}

¶ 2. (a) 在空间 \(\mathbf{R}^n\)（\(n \geqslant 2\)）中，令 U 为补集非空的非空开集；证明 U 的边界含有非空完美集（参见一般拓扑，第一章 § 9，习题 17）。

\begin{enumerate}
\def\labelenumi{(\alph{enumi})}
\setcounter{enumi}{1}
\item
  由 (a) 推出：若 \(\mathbf{A}\) 是 \(\mathbf{R}^n\) 的处处稠密子集，并且与 \(\mathbf{R}^n\) 中每个完美集相交，则 \(\mathbf{A}\) 是凸的。
\item
  证明  中存在一个处处稠密的子域 \(\mathbf{C}\)，它连通且局部连通，并且是 \(\mathbf{K}\) 的纯超越扩张（应用一般拓扑第九章 § 5 的习题 18 (b)、(c)，利用 (b) 以及集合论第三章 § 6，习题 24 中描述的方法，通过超限归纳构造 \(\mathbf{Q}\)）。推出 \(\mathbf{K}\) 中存在包含  的子域 ，它（代数地）同构于 \(\mathbf{K}' \supset \mathbf{K}\)，且连通并局部连通。\(\mathbf{C}\)\(\mathbf{R}\)，且连通并局部连通。
\item
  利用 (c) 证明 C 上存在一种连通且局部连通空间的拓扑，它与域结构相容，并且在此拓扑下 C 的完备化是 C 上的一个代数，且是两个同构于 C 的域的直复合。
\end{enumerate}

\begin{enumerate}
\def\labelenumi{\arabic{enumi}.}
\setcounter{enumi}{2}
\tightlist
\item
  设 K 为全不连通的非离散局部紧交换域；令 A 为 K 上模 \(_{K}\) 的绝对值的环，令 U 为 A 的单位群。对每个整数 n \textgreater{} 0，令 \(_{n}U\) 表示 K 中 n 次单位根所成的群，令 \(U^{n}\) 表示 U 的由 U 中元素的 n 次幂组成的子群。证明若 n 与 K 的特征互素，则 \(U^{n}\) 是 U 的开子群，并且
\end{enumerate}

\[
\operatorname{Card} (\mathrm{U} / \mathrm{U} ^ {n}) = \operatorname{mod} _ {\mathrm{K}} (n). \operatorname{Card} (_ {n} \mathrm{U})
\]

（使用积分论第七章 §2 的习题 14 以及交换代数第三章 §4，第 6 节，定理 2 的推论 1，证明若 m 是 A 的极大理想，则  在  下的像为 \(x \mapsto x^n\)，其中 k 充分大。）\(1 + \mathfrak{m}^k\)\(1 + n \cdot \mathfrak{m}^k\)\(k\)，其中 k 充分大。）

\begin{enumerate}
\def\labelenumi{\arabic{enumi}.}
\setcounter{enumi}{3}
\tightlist
\item
  \begin{enumerate}
  \def\labelenumii{(\alph{enumii})}
  \tightlist
  \item
    设 K 为交换域，v 为 K 上的赋值，且 K 关于 v 为 Hensel（§ 8，习题 6）。进一步假设 K 与 v 的剩余域 k 的特征均为 0。证明 v 的环 A 中存在子域 \(K_{0}\)，使得典范映射 \(A \rightarrow k\) 限制在 \(K_{0}\) 上时，是从 \(K_{0}\) 到 k 的同构。（令 H 为 A 的子域，使 H 在典范映射下的像同构于 k 的子域 E；证明若 \(E \neq k\)，则 A 中存在 \(\alpha \notin H\)，使得 K 的子域 \(\mathrm{H}(\alpha)\) 包含于 A，且典范同构于 \(\mathrm{E}(\bar{\alpha})\)，其中 \(\bar{\alpha}\) 是 \(\alpha\) 在 k 中的类；根据 \(\bar{\alpha}\) 在 E 上代数或超越分两种情形。）
  \end{enumerate}
\end{enumerate}

\begin{enumerate}
\def\labelenumi{(\alph{enumi})}
\setcounter{enumi}{1}
\tightlist
\item
  进一步假设 v 是离散赋值且 K 关于 v 完备；由 (a) 推出 K 同构于形式幂级数域 \(k((\mathbf{T}))\)。
\end{enumerate}

\begin{enumerate}
\def\labelenumi{\arabic{enumi}.}
\setcounter{enumi}{4}
\tightlist
\item
  \begin{enumerate}
  \def\labelenumii{(\alph{enumii})}
  \tightlist
  \item
    设 K 为交换域，v 为 K 上的高度 1 赋值，且 K 关于 v 完备，A 为 v 的环；进一步假设 v 的剩余域 k 是完美域且特征为 p \textgreater{} 0。对每个元素  及每个整数 n，令 \(\xi \in k^{*}\) 为 A 中属于类 \(x_{n}\) 的元素；证明序列 \(\xi^{p-n}\) 是 A 中的 Cauchy 序列，其极限与在类 \((x_{n}^{p^{n}})\) 中选择  的方式无关。若该极限记作 \(x_{n}\)，证明 \(\xi^{p-n}\)\(\phi(\xi)\) 是从乘法群 \(\phi\) 到乘法群 \(k^{*}\) 的唯一同构 u，使得对所有 \(K^{*}\)，\(\xi \in k^{*}\) 是 A 中属于类 \(u(\xi)\) 的元素。\(\xi\) 的元素。
  \end{enumerate}
\end{enumerate}

\begin{enumerate}
\def\labelenumi{(\alph{enumi})}
\setcounter{enumi}{1}
\item
  若  的特征也为 p，证明将 \(\mathbf{K}\) 通过 \(p\)\(\phi\) 延拓到 k 后，它是从域 \(k\)\(\phi(0) = 0\) 到 \(k\) 的一个子域的同构。推出第 3 节定理 1 (iii) 的一个新证明。\(\mathbf{K}\) 的一个子域的同构。推出第 3 节定理 1 (iii) 的一个新证明。
\item
  假设  有限。证明若 \(k\) 与 p 互素，则 \(r\) 中元素的 r 次幂所成群  在 \(p\)\((\mathbf{K}^*)^r\) 中具有有限指数（使用 Hensel 引理）。进一步若 v 离散且 \(r\)\(\mathbf{K}^*\) 的特征为 0，则证明无论 \(\mathbf{K}^*\)\(v\)\(\mathbf{K}\) 如何同样成立（注意 \(r\) 的每个元素都是 p 次幂）。\(1 + p^2\mathbf{A}\)\(p\) 的每个元素都是 p 次幂）。
\end{enumerate}

\exercisesection{10}

\begin{enumerate}
\def\labelenumi{\arabic{enumi}.}
\tightlist
\item
  设 K 为域，P 为 K 的素子域；若 P 的特征为 p \(_{P}\)\textgreater{} 0，则 K 的绝对维数是 dim. al \(_{P}\) K；若 P 的特征为 0，则是 dim. al  K + 1。令 v 为 K 上的赋值，h 为其高度，r 为其有理秩，k 为其剩余域。
\end{enumerate}

\begin{enumerate}
\def\labelenumi{(\alph{enumi})}
\item
  假设 K 的绝对维数  有限。则若 k 的绝对维数为 \(n\)，有 \(K\)\(s\)。\(k\)\(r + s \leqslant n\)。
\item
  进一步假设 K 是 P 的有限生成扩张。则若 \(r + s = n\)，k 是其素子域的有限生成扩张，且 v 的阶群同构于 \(Z'\)；若 \(h + s = n\)，k 是其素子域的有限生成扩张，且 v 的阶群同构于按字典序排序的 \(Z'\)；最后，若 s = n - 1，则 v 是离散赋值，且 k 是其素子域的有限生成扩张。
\end{enumerate}

¶ 2. 设 K 为域，v 为 K 上的赋值，L 为 K 的扩张， 为 L 上延拓 v 的赋值；若 \(v'\)，则称 L（相对于 ）为立即扩张。K 的完备化 \(v'\) 是 K 的立即扩张。\(e(v'/v) = f(v'/v) = 1\)\(\hat{K}\) 是 K 的立即扩张。

\begin{enumerate}
\def\labelenumi{(\alph{enumi})}
\item
  \(\mathbf{L}\) 是 \(\mathbf{K}\) 的立即扩张的充要条件是：对所有 \(x \in \mathbf{L} - \mathbf{K}\)，存在 \(y \in \mathbf{K}\)，使得 \(v'(x - y) > v'(x)\)。
\item
  对每个域 K，证明存在 K 的最大立即扩张 L，即不允许存在不同于自身的立即扩张（使用 § 3，习题 3 (b)）。
\item
  证明若 K 关于由 v 定义的拓扑线性紧，则它除自身外不允许有立即扩张（使用 (a)，并注意在 (a) 的记号中， 时的  所成集合在 \(v'(x - y)\) 的取值群中没有最大元）。\(y \in K\)\(v'\) 的取值群中没有最大元）。
\item
  假设 \(\mathbf{K}\) 不是线性紧的；令 \(\mathbf{B}\) 为 \(v\) 的阶群的一个良序子集，\((a_{\beta})_{\beta \in \mathbb{B}}\) 为 \(\mathbf{K}\) 的一个元素族，满足当 \(\lambda < \mu < \nu\) 时，
\end{enumerate}

\[
v (a _ {\lambda} - a _ {\mu}) <   v (a _ {\mu} - a _ {v}),
\]

但不存在 ，使得对满足 \(x \in \mathbf{K}\) 的每个指标有序对 ，都有 。对 \(v(x - a_{\lambda}) = v(a_{\lambda} - a_{\mu})\) 的每个扩张  以及 v 到  的每个延拓 \((\lambda, \mu)\)\(\lambda < \mu\)\(\mathbf{E}\)，令 \(\mathbf{K}\)\(w\) 表示由满足 \(v\)\(\mathbf{E}\)\(\mathrm{D}_{\mathbf{E}}(\lambda)\) 且 \(z \in \mathbf{E}\) 的 z 所成的集合，其中 \(v(z - a_{\lambda}) \geqslant \gamma_{\lambda}\) 是当 \(\gamma_{\lambda}\) 时  的共同值；因此假设为 \(v(a_{\lambda} - a_{\mu})\)（§ 5，习题 5 (a)）。令 \(\lambda < \mu\)\(\bigcap_{\beta \in \mathbf{B}} \mathrm{D}_{\mathbf{K}}(\beta) = \varnothing\) 为 \(\Omega\) 的代数闭包，令 \(\mathbf{K}\) 为 v 到 \(v_0\) 的延拓；证明：若 \(v\)\(\Omega\) 是 \(\mathbf{P}\) 中的多项式，则存在 \(\mathbf{K}[\mathbf{X}]\) 使得对所有 \(\lambda \in \mathbf{B}\) 有  的充要条件是 \(v(\mathrm{P}(a_{\mu})) = v(\mathrm{P}(a_{\lambda}))\) 在 \(\lambda \leqslant \mu\)\(\mathbf{P}\) 中没有零点属于 \(\Omega\)。若 \(\bigcap_{\beta \in \mathbf{B}} \mathrm{D}_{\Omega}(\beta)\) 具有此性质，且 \(\mathbf{P}\)、w 的含义同上，证明对所有 \(\mathbf{E}\)，当  足够大时，有 \(w\)\(z \in \bigcap_{\beta \in \mathbf{B}} \mathrm{D}_{\mathbf{E}}(\beta), w(\mathrm{P}(z) - \mathrm{P}(a_{\mu})) > w(\mathrm{P}(z)) = v(\mathrm{P}(a_{\mu}))\)\(\mu\)（将  分解为 \(\mathrm{P}(\mathbf{X})\) 中的因子）。推出下列情形之一：\(\Omega[\mathbf{X}]\) 中的因子）。推出下列情形之一：

\begin{enumerate}
\def\labelenumi{(\arabic{enumi})}
\tightlist
\item
  或者 \(\bigcap_{\beta \in B} D_{\Omega}(\beta) \neq \varnothing\)，并且若 \(\theta\) 是该交集中的元素且其在 \(K\) 上的次数尽可能小，则 \(K(\theta)\) 是不同于 \(K\) 的 \(K\) 的立即扩张。
\end{enumerate}

\begin{enumerate}
\def\labelenumi{(\arabic{enumi})}
\setcounter{enumi}{1}
\tightlist
\item
  或者 ；则 K(X) 上存在延拓 v 的赋值 \(\bigcap_{\beta \in B} D_{\Omega}(\beta) = \varnothing\)，使得当 \(v'\) 足够大时 ，并且带有 \(v\)\(v(P(X)) = v(P(a_\mu))\) 的 K(X) 是 K 的立即扩张（使用 (a) 中的判据）。\(\mu\)\(v'\) 的 K(X) 是 K 的立即扩张（使用 (a) 中的判据）。
\end{enumerate}

\begin{enumerate}
\def\labelenumi{(\alph{enumi})}
\setcounter{enumi}{4}
\tightlist
\item
  由 (c) 和 (d) 推出 K 为线性紧的充要条件是 K 除自身外没有立即扩张。
\end{enumerate}

\chapter{除子}

本章所考虑的所有环都假设为交换环并含有单位元。所有环同态都假设将单位元映到单位元。环 A 的每个子环都假设含有 A 的单位元。

\section{KRULL 整环}

\subsection{整环的除性理想}

定义 1. 设 A 为整环，K 为其分式域。K 的每个满足下述条件的 A-子模 a 都称为 A 的分式理想（滥用语言时也称为 K 的分式理想）：存在 A 中元素 \(d \neq 0\)，使得 da ⊂ A。

 的每个有限生成 A-子模  都是分式理想：若 \(\mathfrak{a}\) 是 \(\mathbf{K}\)\((a_{i})_{1\leqslant i\leqslant n}\) 的一组生成元，可写成 \(\mathfrak{a}\)，其中 \(a_{i} = b_{i} / d_{i}\)、\(b_{i}\in \mathbf{A}\) 且 \(d_{i}\in \mathbf{A}\)；若 \(d_{i}\neq 0\)，显然 \(d = d_{1}\ldots d_{n}\)。特别地，\(da\subset A\) 的单生成 A-子模都是分式理想（回忆它们在代数第六章 § 1，第 5 节中被称为分式主理想）。若 \(\mathbf{K}\) 是 Noether 环，每个分式理想都是有限生成的 A-模。\(\mathbf{A}\) 的分式理想的每个 A-子模都是分式理想。\(\mathbf{A}\) 的每个理想都是分式理想；为避免混淆，也将它们称为 \(\mathbf{A}\) 的整理想。\(\mathbf{A}\) 的整理想。

记 I(A) 为 A 的非零分式理想所成的集合。给定 I(A) 的两个元素 a、b，将关系“每个包含 a 的分式主理想也包含 b”记作 （或 \(a \prec b\)）；显然此关系是 I(A) 上的预序。令 R 表示相伴的等价关系“\(b \succ a\) 且 ”（集合论，第三章 § 1，第 2 节），令 D(A) = I(A)/R 为商集；称 D(A) 的元素为 A 的除子，并且对每个分式理想 a  I(A)，将 a 在 D(A) 中的典范像记作 div a（或 div \(\prec b\)\(b \prec a\) a），称 div a 为 a 的除子；若 a = Ax 是分式主理想，则以 \(a \in I(A)\)\(_{A}\) 代替 \(\operatorname{div}(x)\)，并称 \(\operatorname{div}(\mathbf{A}x)\) 为 x 的除子；形如 \(\operatorname{div}(x)\) 的 D(A) 的元素称为主除子。取商后，I(A) 上的预序 \(x\)\(\operatorname{div}(x)\) 在 D(A) 上定义了一个序，记作 \(\prec\)。\(\leqslant\)。

对于每个 ，根据假设存在某个 \(\mathfrak{a} \in \mathrm{I}(\mathbf{A})\) 中的 ，使得 \(d \neq 0\)；因此，包含 \(\mathbf{A}\)\(\mathfrak{a} \subset \mathrm{Ad}^{-1}\) 的分式主理想的交  是 \(\tilde{\mathfrak{a}}\) 的元素。显然，关系 \(\mathfrak{a}\)\(\mathrm{I}(\mathbf{A})\) 等价于关系 \(\mathfrak{a} \prec \mathfrak{b}\)；因此关系 \(\tilde{\mathfrak{a}} \supset \tilde{\mathfrak{b}}\) 蕴含 \(\mathfrak{a} \supset \mathfrak{b}\)。\(\mathfrak{a} \prec \mathfrak{b}\) 的两个元素  模 \(\mathfrak{a}, \mathfrak{b}\) 等价的充要条件是 \(\mathrm{I}(\mathbf{A})\)\(\mathbb{R}\)。\(\tilde{\mathfrak{a}} = \tilde{\mathfrak{b}}\)。

定义 2. I(A) 中每个满足  的元素  称为 A 的除性分式理想。\(\mathfrak{a}\)\(\mathfrak{a} = \tilde{\mathfrak{a}}\) 称为 A 的除性分式理想。

换言之，除性理想恰是非零的非空分式主理想族的交。除性理想的每个非零交仍是除性理想。若 a 是除性理想，则对所有 \(x \in K^{*}\)，ax 也是除性理想，因为映射 \(b \mapsto bx\) 将分式主理想集双射到自身。对所有 \(a \in I(A)\)，\(\tilde{a}\) 是包含 a 的最小除性理想，且与 a 模 R 等价。此外，若 b 是与 a 模 R 等价的除性理想，则 \(\tilde{a} = \tilde{b} = b\)。因此 \(\tilde{a}\) 是满足 div a = div b 的唯一除性理想 b（换言之，将映射 \(a \mapsto div a\) 限制到除性理想集上所得映射是单射）。

设  与 \(\mathfrak{a}\) 为 \(\mathfrak{b}\) 的两个分式理想。回忆（第一章 § 2，第 10 节），\(\mathbf{K}\) : \(\mathfrak{b}\) 表示满足 \(\mathfrak{a}\) 的  所成的集合；这显然是一个 A-模；若 \(x \in \mathbf{K}\) 且 \(x\mathfrak{a} \subset \mathfrak{b}\)\(\mathfrak{b} \in \mathrm{I}(\mathrm{A})\)，则 \(\mathfrak{a} \in \mathrm{I}(\mathrm{A})\)；因为若 \(\mathfrak{b} \colon \mathfrak{a} \in \mathrm{I}(\mathrm{A})\) 是 \(d\) 中满足 \(\mathbf{A}\) 且 \(db \subset \mathbf{A}\) 的非零元素，且 \(da \subset \mathbf{A}\) 是 \(a\) 中的非零元素，则 \(\mathbf{A} \cap \mathfrak{a}\)；另一方面，若 \(da(\mathfrak{b} : \mathfrak{a}) \subset \mathbf{A}\) 属于 \(b \neq 0\)，则 \(\mathfrak{b}\)，故 \(bda \subset \mathfrak{b}\) : \(bd \in \mathfrak{b}\)，且 \(\mathfrak{a}\)。\(\mathfrak{b} : \mathfrak{a} \neq 0\)。

b: a 的定义还可写成：

\[
\mathfrak {b}: \mathfrak {a} = \bigcap_ {x \in \mathfrak {a}, x \neq 0} \mathfrak {b} x ^ {- 1}.\tag{1}
\]

命题 1. (a) 若 \(\mathfrak{b}\) 是除性理想且 \(\mathfrak{a} \in \mathbf{I}(\mathbf{A})\)，则 \(\mathfrak{b}\) : \(\mathfrak{a}\) 是除性理想。

\begin{enumerate}
\def\labelenumi{(\alph{enumi})}
\setcounter{enumi}{1}
\item
令 \(\mathfrak{a},\mathfrak{b}\) 属于 I(A)。要有 \(\operatorname {div}\mathfrak{a} = \operatorname {div}\mathfrak{b}\)，充要条件是 A: \(\mathfrak{a} = \mathbf{A}\) : \(\mathfrak{b}\)。
\item
对所有 \(\mathfrak{a} \in \mathbf{I}(\mathbf{A})\)，都有 \(\tilde{\mathfrak{a}} = \mathbf{A} : (\mathbf{A} : \mathfrak{a})\)。
\end{enumerate}

断言 (a) 由公式 (1) 立即得到，因为若  是除性理想，则对所有 \(\mathfrak{b}\)， 也是除性理想。\(\mathfrak{bx}^{-1}\)\(x \neq 0\) 也是除性理想。

为证明 (b)，令  表示包含 \(\mathrm{P}(\mathfrak{a})\) 的分式主理想所成的集合；关系 \(\mathfrak{a}\) 等价于 \(Ax \in \mathrm{P}(\mathfrak{a})\)，从而等价于 \(x^{-1}\mathfrak{a} \subset A\) : \(x^{-1} \in A\)。由于关系 \(\mathfrak{a}\) 按定义等价于 \(\operatorname{div} \mathfrak{a} = \operatorname{div} \mathfrak{b}\)，它也等价于 A : \(\mathrm{P}(\mathfrak{a}) = \mathrm{P}(\mathfrak{b})\) : \(A\)\(\mathfrak{a} = A\)。\(\mathfrak{b}\)。

最后，由于 ，有 。在此式中以  代替 ，可见 \(\mathfrak{a}(\mathbf{A}:\mathfrak{a})\subset \mathbf{A},\mathfrak{a}\subset \mathbf{A}:(\mathbf{A}:\mathfrak{a})\)\(\mathfrak{a}\)；另一方面，关系 \(\mathbf{A}:\mathfrak{a}\)\(\mathbf{A}:\mathfrak{a}\subset \mathbf{A}:(\mathbf{A}:(\mathbf{A}:\mathfrak{a}))\) 蕴含\(\mathfrak{a}\subset \mathbf{A}:(\mathbf{A}:\mathfrak{a})\) 蕴含

\[
\mathrm{A}: \mathfrak {a} \supset \mathrm{A}: (\mathrm{A}: (\mathrm{A}: \mathfrak {a})).
\]

因此 A: a = A : (A : (A : a))，由 (b) 得 div a = div(A : (A : a))。由于 A : (A : a) 根据 (a) 是除性理想，当然 \(\tilde{a} = A\) : (A : a)，这就证明了 (c)。

注。在上述证明过程中已经证明：对每个理想 \(a \in I(A)\)，都有 A: a = A: (A: (A: a))；这是集合论第三章 § 1，第 5 节命题 2 的特例。

命题 2. (i) 在 D(A) 中，每个非空有上界的集合都有最小上界。更确切地，若 \((\mathfrak{a}_t)\) 是 I(A) 中有上界的非空元素族，则

\[
\sup _ {i} (\operatorname{div} a _ {i}) = \operatorname{div} \left(\bigcap_ {i} \tilde {a} _ {i}\right).
\]

\begin{enumerate}
\def\labelenumi{(\roman{enumi})}
\setcounter{enumi}{1}
\tightlist
\item
在 D(A) 中，每个非空有下界的集合都有最大下界。更确切地，若 \((\mathfrak{a}_{\mathfrak{t}})\) 是 I(A) 中有下界的非空元素族，则
\end{enumerate}

\[
\inf _ {i} (\operatorname{div} a _ {i}) = \operatorname{div} \left(\sum_ {i} a _ {i}\right).
\]

\begin{enumerate}
\def\labelenumi{(\roman{enumi})}
\setcounter{enumi}{2}
\tightlist
\item
集合 D(A) 是一个格。
\end{enumerate}

设 \((\mathfrak{a}_i)\) 是 I(A) 中有上界的非空元素族。说除性理想 \(\mathfrak{b}\) 是这个族的上界，就是说它包含于所有 \(\tilde{\mathfrak{a}}_i\) 中，也就是说 \(\mathfrak{b}\) 包含于 \(\bigcap_{i} \tilde{\mathfrak{a}}_i\) 中。因此 \(\bigcap_{i} \tilde{\mathfrak{a}}_i \neq (0)\)，且 \(\bigcap_{i} \tilde{\mathfrak{a}}_i\) 是除性理想，这就证明了 (i)。

现在令 \((\tilde{a}_t)\) 为 I(A) 中有下界的非空元素族。说除性理想 \(b\) 是这个族的下界，就是说它包含所有 \(\tilde{a}_t\)，也就是说（因为 \(b\) 是除性理想）它包含所有 \(a_t\)，或者说 \(b \supset \sum_{i} a_i\)。这就证明了 (ii)。

最后，为证明 (iii)，根据 (i) 与 (ii)，只需证明当 a、b 属于 I(A) 时，集合 \(\{a, b\}\) 在 I(A) 中既有上界又有下界；现在它以上界 \(a \cap b\) 为界（该集合不同于 (0)）。它以下界 \(a + b\) 为界，因为 \(a + b \in I(A)\)：若 d 与 \(d'\) 是 A 中满足 \(da \subset A\) 且 \(d'b \subset A\) 的非零元素，则 \(dd'(\mathfrak{a} + \mathfrak{b}) \subset \mathbf{A}\)。

推论。若 \(x, y\) 与 \(x + y\) 属于 \(\mathbf{K}^*\)，则 \(\operatorname{div}(x + y) \geqslant \inf(\operatorname{div}(x), \operatorname{div}(y))\)。\(\mathbf{A}(x + y) \subset \mathbf{A}x + \mathbf{A}y\)，从而 \(\operatorname{div}(x + y) \geqslant \operatorname{div}(\mathbf{A}x + \mathbf{A}y)\)。
